%%% Hlavní soubor. Zde se definují základní parametry a odkazuje se na ostatní části. %%%

%% Verze pro jednostranný tisk:
% Okraje: levý 40mm, pravý 25mm, horní a dolní 25mm
% (ale pozor, LaTeX si sám přidává 1in)
\documentclass[12pt,a4paper]{report}
\setlength\textwidth{145mm}
\setlength\textheight{247mm}
\setlength\oddsidemargin{15mm}
\setlength\evensidemargin{15mm}
\setlength\topmargin{0mm}
\setlength\headsep{0mm}
\setlength\headheight{0mm}
% \openright zařídí, aby následující text začínal na pravé straně knihy
\let\openright=\clearpage

%% Pokud tiskneme oboustranně:
% \documentclass[12pt,a4paper,twoside,openright]{report}
% \setlength\textwidth{145mm}
% \setlength\textheight{247mm}
% \setlength\oddsidemargin{15mm}
% \setlength\evensidemargin{0mm}
% \setlength\topmargin{0mm}
% \setlength\headsep{0mm}
% \setlength\headheight{0mm}
% \let\openright=\cleardoublepage

%% Pokud používáte csLaTeX (doporučeno):
\usepackage{czech}
%% Pokud nikoliv:
%\usepackage[czech]{babel}
%\usepackage[T1]{fontenc}

%% Použité kódování znaků: obvykle latin2, cp1250 nebo utf8:
\usepackage[utf8]{inputenc}

%% Ostatní balíčky
\usepackage{graphicx}
\usepackage{amsthm}

\usepackage{parskip}
\setlength{\parindent}{1cm}
\setlength{\parskip}{0.2cm plus 3mm minus 1mm}
\usepackage{indentfirst}
\usepackage{alltt}
\usepackage{listings}


%% Balíček hyperref, kterým jdou vyrábět klikací odkazy v PDF,
%% ale hlavně ho používáme k uložení metadat do PDF (včetně obsahu).
%% POZOR, nezapomeňte vyplnit jméno práce a autora.
%\usepackage[ps2pdf,unicode]{hyperref}   % Musí být za všemi ostatními balíčky
\usepackage[unicode]{hyperref}   % Musí být za všemi ostatními balíčky
\hypersetup{pdftitle=Hybridní databáze}
\hypersetup{pdfauthor=Radek Hušek}

%%% Drobné úpravy stylu

% Tato makra přesvědčují mírně ošklivým trikem LaTeX, aby hlavičky kapitol
% sázel příčetněji a nevynechával nad nimi spoustu místa. Směle ignorujte.
\makeatletter
\def\@makechapterhead#1{
  {\parindent \z@ \raggedright \normalfont
   \Huge\bfseries \thechapter. #1
   \par\nobreak
   \vskip 20\p@
}}
\def\@makeschapterhead#1{
  {\parindent \z@ \raggedright \normalfont
   \Huge\bfseries #1
   \par\nobreak
   \vskip 20\p@
}}
\makeatother

% Toto makro definuje kapitolu, která není očíslovaná, ale je uvedena v obsahu.
\def\chapwithtoc#1{
\chapter*{#1}
\addcontentsline{toc}{chapter}{#1}
}

%\linespread{1.5}
%\parskip 2.2pt

\begin{document}

% Trochu volnější nastavení dělení slov, než je default.
\lefthyphenmin=2
\righthyphenmin=2

%%% Titulní strana práce

\pagestyle{empty}
\begin{center}

\large

Univerzita Karlova v~Praze

\medskip

Matematicko-fyzikální fakulta

\vfill

{\bf\Large BAKALÁŘSKÁ PRÁCE}

\vfill

\centerline{\mbox{\includegraphics[width=60mm]{logo.eps}}}

\vfill
\vspace{5mm}

{\LARGE Radek Hušek}

\vspace{15mm}

% Název práce přesně podle zadání
{\LARGE\bfseries Hybridní databáze}

\vfill

% Název katedry nebo ústavu, kde byla práce oficiálně zadána
% (dle Organizační struktury MFF UK)
Katedra aplikované matematiky

\vfill

\begin{tabular}{rl}

Vedoucí bakalářské práce: & Mgr. Martin Mareš, Ph.D. \\
\noalign{\vspace{2mm}}
Studijní program: & Informatika \\
\noalign{\vspace{2mm}}
Studijní obor: & obecná informatika \\
\end{tabular}

\vfill

% Zde doplňte rok
Praha 2012

\end{center}

\newpage

%%% Následuje vevázaný list -- kopie podepsaného "Zadání bakalářské práce".
%%% Toto zadání NENÍ součástí elektronické verze práce, nescanovat.

%%% Na tomto místě mohou být napsána případná poděkování (vedoucímu práce,
%%% konzultantovi, tomu, kdo zapůjčil software, literaturu apod.)

%\openright

%\noindent
%Poděkování.

%\newpage

%%% Strana s čestným prohlášením k bakalářské práci

\vglue 0pt plus 1fill

\noindent
Prohlašuji, že jsem tuto bakalářskou práci vypracoval samostatně a~výhradně
s~použitím citovaných pramenů, literatury a~dalších odborných zdrojů.

\medskip\noindent
Beru na~vědomí, že se na moji práci vztahují práva a~povinnosti vyplývající
ze zákona č. 121/2000 Sb., autorského zákona v~platném znění, zejména skutečnost,
že Univerzita Karlova v~Praze má právo na~uzavření licenční smlouvy o~užití této
práce jako školního díla podle §60 odst. 1 autorského zákona.

\vspace{10mm}

\hbox{\hbox to 0.5\hsize{%
V~Praze dne ............
\hss}\hbox to 0.5\hsize{%
Podpis autora
\hss}}

\vspace{20mm}
\newpage

%%% Povinná informační strana bakalářské práce

\vbox to 0.5\vsize{
\setlength\parindent{0mm}
\setlength\parskip{5mm}

Název práce:
Hybridní databáze
% přesně dle zadání

Autor:
Radek Hušek

Katedra:  % Případně Ústav:
Katedra aplikované matematiky
% dle Organizační struktury MFF UK

Vedoucí bakalářské práce:
Mgr. Martin Mareš, Ph.D., Katedra aplikované matematiky
% dle Organizační struktury MFF UK, případně plný název pracoviště mimo MFF UK

Abstrakt: Tato práce popisuje návrh a~implementaci datové struktury, která
se snaží kombinovat výhody databází a~běžných datových struktur.
Ze světa databází vychází především podpora pro persistenci dat
prostřednictvím jejich uložení na disku a~práce s~daty pomocí
transakcí, které umožňují paralelní přístup při zajištění
konzistence dat. Od datových struktur naopak přichází implementace
v~podobě knihovny funkcí a~snaha o~maximální jednoduchost
a~uložení dat v~paměti.
Navržená databáze staví na konceptu transakční paměti a~data jsou
na disku ukládána ve formě záznamu provedených operací.

% abstrakt v rozsahu 80-200 slov; nejedná se však o opis zadání bakalářské práce



Klíčová slova: in-memory databáze, transakční paměť, datové struktury
% 3 až 5 klíčových slov

\vss}\nobreak\vbox to 0.49\vsize{
\setlength\parindent{0mm}
\setlength\parskip{5mm}

Title:
Hybrid Databases
% přesný překlad názvu práce v angličtině

Author:
Radek Hušek

Department:
Department of Applied Mathematics
% dle Organizační struktury MFF UK v angličtině

Supervisor:
Mgr. Martin Mareš Ph.D., Department of Applied Mathematics
% dle Organizméno a příjmační struktury MFF UK, případně plný název pracoviště
% mimo MFF UK v angličtině

Abstract: This thesis presents design and implementation of a~data structure, which tries
to combine advantages of both databases and regular data structures.
Main advantages of databases we try to retain are data persistence
through storing data on a~hard disk and working with data using
transactions which allows us parallel access
without danger of inconsistency. From data structures
we borrow the implementation as a~library of functions and the aim on simplicity
and storing data in memory.
Our implementation is built around the concept of (software) transactional
memory; all data are stored on hard drive as log of operations.


% abstrakt v rozsahu 80-200 slov v angličtině; nejedná se však o překlad
% zadání bakalářské práce



Keywords: in-memory database, transactional memory, data structure
% 3 až 5 klíčových slov v angličtině

\vss}

\newpage

%%% Strana s automaticky generovaným obsahem bakalářské práce. U matematických
%%% prací je přípustné, aby seznam tabulek a zkratek, existují-li, byl umístěn
%%% na začátku práce, místo na jejím konci.

\openright
\pagestyle{plain}
\setcounter{page}{1}
\tableofcontents

%%% Jednotlivé kapitoly práce jsou pro přehlednost uloženy v samostatných souborech
\chapter*{Úvod}
\addcontentsline{toc}{chapter}{Úvod}

Tato práce se snaží vyplnit mezeru, která se rozevírá mezi (relačními) databázovými systémy
na straně jedné a~klasickými datovými strukturami na straně druhé. Dnešní
 databáze se zaměřují na zpracování obrovských množství dat, a~proto
užívají velmi sofistikované metody manipulace s~nimi. Tyto
pokročilé postupy mají však i~svou stinnou stránku -- nemalou časovou i~prostorovou
režii. Ta se u~obrovských souborů dat neprojevuje,
neboť je bohatě kompenzována vyšší efektivitou zpracování,
ale u~malých dat bývá významná.
Oproti tomu datové struktury jsou velmi rychlé a~mají flexibilnější rozhraní,
které lze často přizpůsobit konkrétní aplikaci, ale postrádají možnost perzistentně
ukládat data (tedy tak aby přežila restart či pád aplikace).

V~důsledku zvětšování velikosti operační paměti počítačů
stále častěji nastává situace, kdy se potřebná data vejdou do operační paměti.
Pokud je u~dat požadováno trvalé uložení mezi běhy aplikace a~jsou tak cenná,
že je nutné zajistit možnost jejich obnovy i~v~případě havárie systému,
nezbývá programátorovi, než místo rychlé datové struktury použít databázi.
Použití databáze zajistí potřebnou perzistenci a~konzistentnost dat, ale
oproti použití datové struktury negativně ovlivní výkon aplikace.

Z~toho důvodu jsem se rozhodl vytvořit hybridní databázi, která se snaží kombinovat
výhody databází a~datových struktur. Udržuje všechna data v~paměti, 
ve tvaru co nejpřátelštějším k~programovacímu
jazyku dané aplikace. Zajišťuje uložení jejich konzistentní verze na disku
a~umožňuje pracovat s~nimi pomocí transakcí. Podpora transakcí je důležitá
jak z~důvodu určení konzistentní verze dat pro uložení na disku, tak především
s~ohledem na stále narůstající počet (logických) procesorů v~dnešních
počítačích, neboť elegantně eliminuje potřebu složité ruční synchronizace přístupu
k~datům,
která je velmi častým zdrojem těžko odhalitelných chyb.

Součástí práce je i~ukázková implementace navržené databáze v~jazyce C ve~formě
knihovny funkcí. Jazyk C byl zvolen pro svou rozšířenost, a~to především v~unixovém
světě. Výhodou jazyka C oproti C{\tt++} je jeho rozšířenost a~\uv{kompatibilita} -- zatímco
užívat z~C{\tt++} knihovny vytvořené v~C je velmi snadné, je použití knihoven
vytvořených v~C{\tt++}, především obsahuje-li jejich rozhraní šablony, v~podstatě nemožné
(přičemž podobná situace nastává s~vazbami do jiných programovacích jazyků).

\subsection*{Rozčlenění práce}
\begin{itemize}
  \item První kapitola obsahuje motivaci k~vytvoření hybridní databáze a~hrubý
    popis jejího fungování.
  \item Druhá kapitola se zabývá softwarovou transakční pamětí a~její implementací
    použitou v~této práci.
  \item Třetí kapitola popisuje formát uložení databáze na disku.
  \item Čtvrtá kapitola obsahuje popis rozhraní navržené databáze.
  \item V~páté kapitole jsou rozebrány implementačně závislé aspekty hybridní
    databáze.
  \item Příloha 1 obsahuje příklad jednoduchého programu, který využívá hybridní databázi.
  \item V~příloze 2 naleznete výsledky testů škálovatelnosti.
\end{itemize}



\chapter{Popis hybridní databáze}
V~této kapitole stručně popíšeme z~jakých předpokladů jsme při návrhu hybridní
databáze vycházeli, jak se promítly do požadavků na ni kladených a~jak se
nám tyto povedlo naplnit. Závěr kapitoly je věnován popisu základního
fungovaní navržené struktury, který je hlouběji rozveden v~kapitolách následujících.

\section{Motivace}
Jak bylo zmíněno v~úvodu, hlavním motivem této práce je nabídnout alternativu
programům, které potřebují perzistentně uchovávat data, a~přitom nepotřebují
komplexní databázový systém. Pro začátek uveďme, jaké klíčové vlastnosti
od vytvářené databáze chceme:

\begin{itemize}
  \item Implementovat databázi jako knihovnu funkcí, takže komunikace
    s~databází znamená přímé volání těchto funkcí bez dalších mezivrstev.
  \item Ukládat data v~paměti, takže čtení je vždy rychlé, a~přitom paralelně
    udržovat obraz dat na disku, aby v~případě havárie programu bylo možné
    obnovit data v~konzistentní verzi.
  \item Vyhledávání dat místo dotazovacího jazyka zajišťují uživatelem vytvořené
    indexy. Ty mají s~indexy v~relačních databázích společné, že obsah obou je
    jednoznačně určen obsahem databáze. Na rozdíl od nich je jejich rozhraní mnohem
    flexibilnější -- uživatel komunikuje s~indexem přímo a~rozhraní indexu
    určuje jeho tvůrce, není pevně zadrátováno v~implementaci databáze.
  \item Podpora více vláken a~udržování konzistence dat za pomoci transakcí.
  \item Skládání operací do větších celků pomocí zanořování transakcí.
\end{itemize}

\subsection{Knihovna funkcí vs. samostatný proces}
Většina databázových systémů, a~to i~nerelačních, se drží koncepce samostatného
databázového serveru, s~nímž klient komunikuje (typicky protokolem TCP/IP).
To je výhodné pro velké systémy, neboť je možné server samostatně ovládat
a~přistupovat k~němu z~více aplikací. Navíc získáme síťovou transparentnost a~větší
spolehlivost, protože pád aplikace neublíží serveru a~aplikace také nevidí
paměť databáze, takže ji nemohou poškodit.

Pro naše potřeby, kdy počítáme s~aplikací menších rozměrů, je síťová
transparentnost nadbytečná a~samostatné ovládání serveru spolu s~přístupem
k~němu z~více aplikací také není předpokládaným použitím. (Ačkoliv
s~úpravami ohledně alokace paměti je teoreticky možné navrhovanou
databázi provozovat jako samostatný proces, který komunikuje s~ostatními
přes sdílenou paměť.)

Naopak díky integraci do klientského procesu získáváme rychlou komunikaci
s~databází spočívající v~přímém volání jejích metod. A~jako
identifikátory jednotlivých objektů uložených v~databázi můžeme používat přímo
ukazatele, což opět výrazně snižuje režii.

\subsection{Primární úložiště v~paměti, záloha na disku}
Jelikož cílíme na menší objemy dat, je výhodné implementovat hybridní
databázi, aby pracovala s~daty přímo v~paměti, a~jejich kopii na disku užívat
jen jako zálohu,
ze které budou obnovena při příštím startu. Tím sice zvýšíme paměťové
nároky (dat nicméně není tolik, aby to představovalo výrazný problém), ale
zároveň nemusíme mít diskový formát, ve kterém lze rychle vyhledávat.
Můžeme tedy upustit od formátů, které respektují strukturu dat,
a~databázi zapisovat na disk jako log -- tedy posloupnost operací, které
vedly k~jejímu současnému stavu.
(Log samozřejmě nemusí obsahovat všechny události od vytvoření databáze. Běžný
postup je, že pokud velikost logu přeroste rozumnou mez, je vygenerován nový,
který obsahuje nejkratší posloupnost operací vedoucí k~současnému stavu.)

To výrazně snižuje jak složitost
zápisu, tak i~počet I/O operací nutných k~jeho provedení (protože
nezapisujeme na přeskáčku, ale vždy jen na konec souboru), což
se významně projeví především u~rotačních pevných disků (ačkoli i~pro
disky SSD je tento přístup výhodnější, neboť jednotlivé zápisy jsou výrazně
větší). Navíc při použití datového formátu respektujícího strukturu dat
by stejně bylo nutné udržovat logu nebo jiným způsobem zajistit atomicitu
prováděných operací.

\subsection{Organizace dat v~paměti}
Data v~paměti budeme chceme organizovat podobně jako datové struktury. Základní
jednotkou dat je {\bf uzel}. Každý uzel se skládá z~{\bf atributů}, které obsahují
vlastní data, přičemž jak atributy jaké atributy obsahuje je určeno jeho typem.
Při porovnání s~relační databází typy uzlů odpovídají tabulkám, uzly řádkům tabulek
a~atributy jednotlivým buňkám. Z~hlediska objektového programování, které je bližší
skutečné implementaci, jsou typy uzlů třídami, uzly instancemi těchto
tříd a~atributy jednotlivými instančními proměnnými.

Objektovému pohledu také odpovídá, že každý uzel má svou jednoznačnou identitu
-- dokážeme rozlišit dva uzly, i~když mají stejný obsah. Za běhu je vhodné jako
identifikátor uzlu používat ukazatel na něj.

Atributy nemohou obsahovat libovolné hodnoty, ale každý má pevně určený typ. Typy
atributů pokrývají základní typy jazyka C -- jsou mezi nimi tedy celá čísla
různých rozsahů s~či bez znaménka a~čísla s~plovoucí řádovou čárkou -- a~obsahují
několik užitečných typů navíc -- řetězec schopný pojmout libovolnou posloupnost
bytů (který se ale chová jako jednoduchý hodnotový typ, ne jako pole) a~ukazatel na jiný
uzel.

\subsection{Uživatelsky definované indexy}
Na rozdíl od relačních databází, kde indexy slouží \uv{pouze} k~urychlení
dotazů, v~navrhované hybridní databázi hrají, především kvůli absenci
dotazovacího jazyka, klíčovou roli. Indexem rozumíme uživatelem definovanou
datovou strukturu jejíž stav (myšlen je samozřejmě stav z~pohledu vnějšího
světa, ne zakódování tohoto stavu uvnitř indexu) je jednoznačně určen stavem
databáze bez ohledu
na posloupnost operací, které k~jejímu současnému stavu vedly.

Z~pohledu databáze je index určen callback funkcí, kterou databáze volá při
vytvoření, smazaní či změně uzlu, a~kontextu, který je předán callback funkci.
Uživatelské rozhraní indexu určuje jeho tvůrce
s~tím, že by uživatelské
funkce neměly index modifikovat (protože jeho stav má být určen pouze stavem
databáze). U~každého indexu je při jeho definici možno určit, zda bude
sledovat všechny uzly v~databázi, nebo bude vázán pouze na určitý typ uzlu.
Tento způsob svázání indexů s~databází byl zvolen, protože představuje rozumný
kompromis.

Pokud by všechny indexy byly globální a~uživatel by filtroval zajímavé
události až v~callback funkci, mohlo by to znamenat nemalou režii, protože
většina volání by končila zjištěním, že daný index událost nezajímá.
Na druhou stranu
možnost mít globální index je nezbytná, protože není možné ho přímočaře nahradit
více menšími indexy. Jedinou zbývající možností je nechat uživatele při registraci
indexu určit, k~jakým uzlům je vázán. 

Zde se nabízí mnoho variant, jak takovou volbu
provést. Například vytvořit dotazovací jazyk, který by uživateli dovolil filtrovat
události podle typu uzlu a~jeho obsahu. 
Tím by bylo možné vyhnout se zbytečným voláním callback funkcí, ovšem režie
takto komplexního systému by nejspíše převážila nad získanou úsporou.
Proto jsem zvolil co nejjednodušší metodu, která by zároveň měla vyhovovat většině
nasazení -- tedy že indexy jsou registrované jako globální, nebo jsou vázány na
jeden typ uzlu.

\subsection{Podpora transakcí}
Podpora transakcí je klíčová ze dvou důvodů. Za prvé databáze musí mít přehled,
kdy jsou data v~ní obsažená v~konzistentním stavu, aby je jako taková označila
při zápisu na disk. Kdyby tak nečinila, bylo by obtížné po pádu aplikace zjistit,
zda jsou data poškozena, a~nemožné obnovit korektní stav dat bez složité
spolupráce uživatele databáze. Druhým a~neméně významným důvodem pro používaní
transakcí je velké zjednodušení práce s~databází z~více vláken současně, jelikož
odpadá komplikovaná ruční synchronizace.


\section{Abstraktní popis rozhraní databáze}
Jak bylo řečeno v~úvodu, hybridní databáze bude implementována jako knihovna funkcí
v~jazyce C. Protože definovat typy uzlů, index a~databází v~jazyce C by bylo velmi
nepohodlné, bude k~tomuto účelu vytvořen samostatný jazyk, který bude překládán
do jazyka C. Popis jazyka definic i~skutečného rozhraní databáze naleznete 
ve čtvrté kapitole.
Zde pouze zmíním základní principy fungování databáze.

\subsection{Transakce}
Transakce budeme implementovat pomocí konceptu softwarové transakční paměti (STM).
Každé vlákno, které bude chtít s~databází pracovat, si nejprve vytvoří {\bf handle}
-- objekt skrze, který bude probíhat další práce. Kdybychom nezavedli handle,
musela by každá transakce existovat jako objekt. To by znamenalo, že každá transakce
na počátku musela provést alokaci tohoto objektu a~na svém konci jeho uvolnění.
To je možné implementovat celkem efektivně, ale databáze navíc potřebuje
vědět o~bežících transakcích. Proto by se tyto musely vždy nějak registrovat,
a~tedy modifikovat hlavní strukturu popisující databázi, což není žádoucí.

Pokud máme handly, tak každá transakce existuje uvnitř příslušného handlu. Ten je
lokální pro dané vlákno, takže není problém s~přístupem k~němu. Aby měla databáze
přehled o~probíhajících transakcích, musejí se i~handly registrovat. V~případě handlů
to ale nevadí, protože na rozdíl od transakcí existují dlouho a~vznikají málokdy.
(Typický scénář je, že každé vlákno si vytvoří jeden handle, který bude používat po
celou dobu svého běhu.)
Nad handlem jsou definovány všechny ostatní funkce pro práci s~databází. Mezi základní
operace patří začátek a~konec transakce, zrušení probíhající transakce a~funkce
pro práci s~uzly (tedy vytvoření a~smazání uzlu, přečtení hodnoty atributu
a~zápis hodnoty atributu). 

\subsection{Indexy}
Indexy byli již popsány výše v~této kapitole, takže pouze pár doplnění.
Jaké indexy databáze obsahuje je vždy definováno už před jejím vytvořením
a~za běhu indexy není možné přidávat ani odebírat. Naopak v~diskovém
formátu se přítomnost indexů nijak neprojeví -- je tedy možné používat
tutéž databázi v~různých aplikacích s~různými indexy.

Všem indexům je společné,
že z~pohledu databáze jsou tvořeny kontextem a~callback funkcí. Kontexty
indexů jsou součástí transakční paměti. 
To je nezbytné, aby stav indexů zůstal konzistentní se stavem databáze.
Kvůli tomu je i~s~indexy nutné pracovat výhradně pomocí transakcí.
O~alokaci kontextů se databáze stará sama a~uživateli nejsou přímo přístupné.
S~tím souvisí, že z~pohledu uživatele jsou metody indexy obaleny makrem,
které se stará o~předání správného kontextu.


\chapter{Softwarová transakční paměť}

\section{STM obecně}

Softwarová transakční paměť je koncept pro řízení přístupu ke sdíleným datům.
Oproti řešením založeným na zámcích je to koncept optimistický -- tedy zatímco
se zámky vlákno vždy nejprve musí získat potřebné zdroje (zámky) a~až poté
provést svou činnost, při použití STM vlákno začne optimisticky vykonávat, co
potřebuje, prostředky získává implicitně za běhu, když je potřebuje, a~případný
konflikt s~jiným vláknem se řeší, až k~němu dojde. Typicky se konflikt řeší tím,
že jedno z~vláken odstraní provedené změny, čímž uvolní prostředky, takže druhé
může pokračovat, a~zrušené vlákno se zkusí provést později.


\subsection{Fungování transakcí}
Základním kamenem transakční paměti jsou transakce. Transakce je posloupnost
operací se sdílenými daty, která se vnějšímu pozorovateli jeví atomicky. Tedy
každé vlákno kromě toho, které danou transakci provádí, vidí sdílená data tak,
že transakce buď nebyla provedena vůbec, nebo byla provedena celá, ale nikdy
nepozoruje žádné mezistavy. Z~vnějšího pohledu transakce převádějí data
z~jednoho konzistentního stavu do jiného konzistentního stavu.

Protože jsou transakce vykonávány optimisticky a~prostředky pro svůj běh získávají,
až když je potřebují, může se stát, že požadované prostředky již vlastní jiná transakce.
Takováto situace se nazývá {\bf kolize}. Řešit kolizi čekáním, než jsou potřebné zdroje
uvolněny, není obecně možné, protože není globálně dáno žádné pořadí, v~němž jsou
zdroje získávány. Kvůli může snadno nastat deadlock (tedy situace, kdy graf vztahu
\uv{čekám až transakce $X$ uvolní zdroje, které potřebuji} obsahuje orientovaný
cyklus). Standardním řešením je v~případě kolize zrušit změny provedené jednou
z~transakcí -- provést tzv. {\bf rollback} -- čímž se uvolní prostředky a~druhá
transakce může pokračovat. 

Aby bylo možné transakci zrušit, musí provádět pouze
operace, které je možné vrátit zpět (typicky může zapisovat do paměti, ale nesmí
provádět vstup a~výstup či komunikaci s~jinými procesy). Všechny všechny operace
vykonané v~průběhu transakce jsou zaznamenávány do tzv. {\bf logu}. Rozlišujeme
dva základní typy logů -- zpětný a~dopředný. {\bf Zpětný (undo) log} zaznamenává
původní hodnoty dat, která transakce přepisuje. {\bf Dopředný (redo) log} naopak
funguje tak, že požadavky na zápis nejsou provedeny ihned, ale jsou uloženy do logu
a~provedeny až na konci transakce, když už víme, že kolize nenastala.

Vraťme se k~restartování transakcí.
Pokud budeme -- tak jako v~této práci -- vždy rušit transakci, která chce získat
nedostupné zdroje, vystavujeme se při vysokém zatížení systému riziku livelocku
(situaci, kdy se transakce
vzájemně ruší, aby získaly zdroje, ale žádná není dokončena, protože dříve, než se
stihne dokončit je zrušena). Livelocku se dá vzdorovat zavedením nějaké varianty priorit
transakcí, kdy několikrát restartovaná transakce získá vyšší prioritu, a~při kolizi
je zrušena transakce s~nižší prioritou, nebo snižováním limitu současně běžících
transakcí, pokud nastává příliš mnoho kolizí. Za variantu snižování počtu současně
běžících transakcí lze považovat i~exponenciálně dlouhé čekání před restartem transakce.

V~závislosti na implementaci může transakce detekovat kolizi v~okamžiku, kdy
k~ní dojde, nebo až při pokusu o~úspěšné ukončení transakce, tzv. {\bf commitu}.
Kolize při commitu nejčastěji nastává v~případě, kdy zjistíme, že data čtená
transakcí (tzv. {\bf readset}) byla změněna jinou transakcí.


\subsection{Historie a~budoucnost transakční paměti}
Poprvé byla transakční paměť představena v~článku \cite{HTM}, kde byla předvedena
hardwarová transakční paměť -- počítala tedy s~podporou ze strany hardware. Jako
taková se příliš neujala, ale v~roce 1995 byl i~na základě jejích myšlenek publikován
článek \cite{STM95}, kde byla představena softwarová varianta transakční paměti,
která nevyžadovala (explicitní) hardwarovou podporu. Na tuto práci posléze
navázalo mnoho dalších.

Za pozornost také stojí snaha integrovat podporu transakční paměti do překladače
GCC \cite{gccTM} a~hardwarová implementace transakční paměti, která bude obsažena
v~procesorech generace Haswell firmy Intel \cite{haswellTM}. 
Integrace do překladače a~s~tím spojené rozšíření jazyka
o~atomické konstrukce umožní snadnější vývoj a~hardwarová podpora může přinést výrazné
vylepšení výkonu. Jelikož ale implementace transakční paměti v~GCC je zatím ve velmi
raném stádiu a~procesory Haswell budou uvedeny na trh až v~příštím roce, nebudu
se jim zde dále věnovat.


\section{STM TL2 engine a~jeho rozšíření}

V~této práci jsem jako ideový základ použil
rozšířenou verzi TL2 enginu z~dizertační práce \cite{STMcur}. TL2 engine
vznikl v~roce 2006 \cite{TL2engine}. V~původní implementaci pracoval s~redo logem,
zamykal paměť po slovech a~zámky tedy byly uloženy v~globální tabulce. V~roce
2007 v~práci \cite{STMcur} byl tento engine rozšířen o~podporu undo logu, zamykání
objektů a~umísťování zámků do objektů.

Kromě toho byla implementace výkonnostně optimalizována např. zjednodušením operace
čtení, kde bylo nahrazeno dvojí kontrolování zámku (před a~po načtení hodnoty)
jednou kontrolou po načtení.

V~následujících sekcích
jsou probrány jednotlivé aspekty rozšířeného TL2 enginu
spolu s~popisem a~zdůvodněním postupů zvolených v~implementaci hybridní databáze.


\subsection{Průběh transakce}

\subsubsection*{Globální hodiny a~verzované zámky}
Základním prvkem STM stavící na TL2 enginu jsou verzované zámky spolu s~globálními
hodinami. {\bf Verzovaný zámek} je neblokující zámek,
který ve stavu odemčeno uchovává i~informaci
o~době poslední změny a~v~uzamčeném stavu obsahuje ukazatel na transakci, která
ho uzamkla. Neblokující zámek znamená, že jím lze chránit data proti paralelnímu
přístupu, ale operace \uv{zamkni} nikdy neblokuje -- odpovídá tedy operaci
\verb|trylock| běžných zámků.
Typicky je implementován pomocí 8-bajtového slova, jehož nejnižší
bit určuje stav zamčeno/odemčeno a~zbylé bity podle situace ukazatel na vlastníka 
či verzi
(používáme-li zarovnané ukazatele, nevadí že nemůžeme uložit jejich nejnižší bit,
neboť je vždy nulový). 

{\bf Globální hodiny} jsou vlastně verzovaný zámek, který se vždy
nachází v~odemčeném
stavu a~zvětšuje svou hodnotu s~ukončením (i~neúspěšným) každé transakce,
která zapisovala.

Transakce začíná načtením hodnoty globálních hodin.  Při zápisu (skutečném zápisu,
který v~případě dopředného logu nastává až při commitu) je nejprve
získán příslušný zámek
(při tom se kontroluje, má-li hodnotu nejvýše rovnou začátku transakce),
následně je hodnota zapsána. Zámky jsou uvolněny na konci transakce.

Každé čtení nejprve
zkopíruje čtenou hodnotu do privátní paměti (tj. paměti přístupné jen danému vláknu)
a~poté zkontroluje, zda příslušný
zámek je odemčený a~jeho hodnota je menší nebo rovna času začátku transakce (nebo může
být zámek zamčen, je-li vlastněn danou transakcí).

 Teoreticky by stačilo kontrolovat
readset až při commitu, čímž by se snížila režie. Vznikl by však problém,
že transakce může vidět nekonzistentní data.
U~takovéto transakce by sice commit
nikdy neuspěl, takže by konzistence databáze nebyla porušena, ale v~kódu transakce
by přestaly platit invarianty, které mají o~datech platit. Pokud s~tím kód nepočítá
(a~může být obtížné a~pracné ošetřit všechny možné patologické případy),
může se program např. dostat do nekonečné smyčky. Nebo jsou-li načtená data ukazatel,
může nastat neoprávněný přístup do paměti, protože není zaručeno, kam ukazuje.

Kvůli výkonově kritickým částem je v~implementaci hybridní databáze možnost
kontrolu čtených dat obejít přímým přístupem k~položkám uzlu. Po načtení dat
je však nutné použité uzly explicitně ověřit, aby byly zařazeny do readsetu. Také
je nezbytné ověřovat stav všech ukazatelů dříve než jsou dereferencovány.


\subsubsection*{Schéma průběhu transakce}

Průběh transakce je následující:
\begin{enumerate}
  \item Načtení globálních hodin.
  \item Vlastní obsah transakce -- čtení a~zápisy dat.
  \item V~případě redo logu jsou zabrány všechny zámky potřebné k~zápisu.
  \item Je zkontrolován readset -- tedy že všechny zámky, jimiž chráněná data byla
        čtena, mají hodnotu nejvýše počátku transakce či jsou aktuální transakcí
        vlastněny.
  \item V~případě redo logu jsou všechna data zapsána.
  \item Globální hodiny jsou inkrementovány.
  \item Jsou uvolněny zámky. Jejich nová hodnota odpovídá nové hodnotě globálních
        hodin.
\end{enumerate}

V~případě, že transakce pouze četla a~čtená data průběžně validovala,
je možné kroky 3 až 7 přeskočit. Transakce
totiž nedrží žádné zámky, a~nemusí tedy modifikovat globální hodiny. Kontrolu
readsetu lze pominout, jelikož i~kdyby případná validace skončila neúspěchem,
nemůže tím být poškozena integrita databáze (protože ta nebyla transakcí měněna)
a~okolní svět nemůže poznat, v~jakém pořadí transakce skutečně provedly commit.

\subsection{Zamykání bloků vs. zamykání objektů}

Důležitým rozhodnutím pro výkon STM je volba vhodného modelu zamykání paměti.
Zamykat paměť lze po blocích -- typicky přirozeně zarovnaných slovech či
řádcích cache paměti
dané architektury -- či lze zamykat jednotlivé objekty (ve smyslu úseků paměti,
které mají nějaký rozumný význam). Výhodou zamykaní po blocích je, že STM engine
nemusí nic vědět o~datech a~kromě využívání funkcí STM ke čtení a~zápisu
dat není potřeba jiná spolupráce. Naopak problémem může být zápis do paměti
přes hranice bloků či nevhodná granularita -- zamykání po slovech může být příliš
jemné, zatímco u~zamykání po řádcích cache paměti hrozí, že tentýž zámek hlídá zcela
nesouvisející data. Navíc zámky musejí být uloženy v~globální tabulce
a~jsou blokům přiřazeny pomocí hashovací funkce -- a~zde opět hrozí
riziko, že jeden zámek hlídá nesouvisející data.

Avšak u~tabulky zámků určuje kolize hashovací
funkce a~přesné adresy dat v~paměti, a~tedy mají více náhodný charakter.
Jejich pravděpodobnost lze také snížit zvětšením tabulky, což může mít
negativní dopad na cache paměti (tabulka v~nich zabírá více místa) a~na rychlost
transakcí pokud je např. readset ukládán jako bitmapa tabulky zámků.

Naopak zamykání objektů vyžaduje spolupráci se STM enginem, který musí být
schopen každé čtené či zapisované adrese přiřadit objekt (třeba tak, že rozhraní bude
\uv{čti položku z~objektu} místo \uv{čti z~adresy}). Dále je možné zámky ukládat i~přímo
do objektů místo do globální tabulky, takže nemohou nastávat problémy, kdy jeden
zámek chrání nesouvisející data. To samozřejmě zvýší paměťovou náročnost, jelikož
zámků globální tabulce bývá podstatně méně než objektů.

\subsubsection*{Zvolená implementace}
Z~pohledu granularity je, po vzoru mnoha relačních databázových enginů,
nejvhodnější zamykání po uzlech jakožto nejmenších samostatných objektech.
Stejně jsou zamykány indexy, ačkoli u~nich je až konkrétní implementací určeno,
co je objekt (a~tedy kolika zámky je celý index chráněn).

Původním cílem bylo implementovat zamykání s~globální tabulkou, protože v~běžných
implementacích STM výkonem nezaostává za zámky v~objektech a~přitom je implementačně
jednodušší a~šetří paměť. Po provedení měření jsem implementoval i~zamykání po objektech,
jelikož implementace s~globální tabulkou zámků špatně škálovala v~závislosti na počtu
vláken. Tato odchylka od běžných STM je způsobena především prací navíc
při commitu -- tedy zpracováním transakčního logu, aby ho bylo možné zapsat
na disk.

V~případě, že jsou transakce velmi malé (a~řešení
s~globální tabulkou tedy v~průběhu transakce může škálovat), je jich velký počet,
takže čas strávený při commitu je významný (a~nebo je zahlceno servisní vlákno,
navíc vysoké množství commitů přispívá ke cachovému \uv{ping pongu} mezi transakčními
vlákny a~servisním vláknem). A~jsou-li transakce větší, tak se pravděpodobnost kolize
velmi rychle blíží jedné, takže se transakce vzájemně serializují. Použití
globální tabulky větší než několik stovek zámků také není vhodné, protože
příliš velká tabulka zabírá příliš místa v~cache paměti, a~tím výrazně zpomaluje program.


\subsection{Zpětný vs. dopředný log}

Jelikož transakce musí být v~případě kolize schopny odčinit provedené změny,
je nezbytné nějakým způsobem tyto změny zaznamenávat. K~tomu existují dva přístupy
-- zpětný (undo)  a~dopředný (redo) log.
Při použití zpětného logu transakce zapisuje
do logu hodnoty sdílených dat, které změnila, a~při kolizi projde log
pozpátku a~obnoví přepsaná data. Naopak dopředný log zaznamenává požadavky na
zápis do sdílené paměti a~skutečně je provede až při commitu, poté co bylo
ověřeno, že žádná kolize nenastává.

\subsubsection*{Zpětný log}
Výhodou zpětného logu je jeho jednoduchost, kdy k~implementaci stačí prostý
zásobník, a~minimum práce navíc při commitu, kdy po jeho úspěšném provedení
stačí jen uvolnit paměť zpětného logu. Průběh zápisu se zpětným  logem je
přímočarý.
Nejprve získáme potřebný zámek, původní hodnotu přidáme na zásobník
(spolu s~informacemi o~její adrese) a~zapíšeme novou hodnotu. Nastane-li
rollback transakce, projdeme zásobník od vrcholu a~obnovíme hodnoty, poté
odemkneme zámky.

Nevýhodou této implementace je, že velikost zásobníku je přímo úměrná počtu zápisů,
ačkoli míst, kam bylo zapisováno, může být mnohem méně. Za jistých velmi specifických
okolností (mezi zámky a~objekty je bijekce, do logu ukládáme vždy celý objekt
a~nemáme vnořené transakce) lze prostorovou složitost této implementace snížit
na lineární s~počtem modifikovaných míst tím, že do logu zapíšeme hodnotu,
pouze pokud daný zámek opravdu zamkneme, a~ne pokud ho již vlastníme.

Pokud potřebujeme optimální prostorovou složitost a~se zásobníkovou implementací
jich nelze dosáhnout, nezbývá než užít k~ukládání logu datovou strukturu jako
strom či hashovací tabulka, což zvýší režii, v~případě stromu dokonce logaritmus-krát
zpomalí každou transakci.

\textbf{Poznámka:} Protože uvolnění paměti obecně nelze vrátit, je vhodné paměť,
která měla
být uvolněna v~průběhu transakce, zapisovat do logu a~uvolnit ji až při úspěšném
commitu. (Tedy z~hlediska uvolňování paměti fungovat jako dopředný log.)

\subsubsection*{Dopředný log}
Dopředný log je na implementaci složitější, ovšem oproti zpětnému logu minimalizuje
dobu držení zámků, a~tím usnadňuje paralelní běh transakcí. Větší složitost je
způsobena nutností při každém čtení kontrolovat, zda čtená data nebyla modifikována
a~nejsou tedy uložena v~logu. Kvůli tomu musí být modifikovaná data
uložena v~hashovací tabulce či podobné struktuře s~možností rychlého vyhledávání
a~modifikace.

Výjimkou může být situace, kdy víme, že transakce sebou modifikovaná
data nepotřebuje číst, či jí nevadí číst jejich starou verzi.
(Což zjevně neplatí, chceme-li podporovat skládání transakcí.)
V~takovémto případě lze dopředný log
implementovat pomocí fronty, jejíž každý prvek je požadavek na zápis.
Další nevýhodou může být prodloužení commitu, neboť v~jeho průběhu je nutné
veškerá modifikovaná data skutečně zapsat do sdílené paměti.

\textbf{Poznámka:} Alokaci paměti není samozřejmě možné odsunout a~musí být alokována
již v~průběhu transakce, takže se s~ní musí pracovat pomocí zpětného logu.

\subsubsection*{Zvolená implementace}
Transakce typicky více čtou než zapisují,
takže optimalizace čtení má větší efekt než optimalizace zápisu.
Proto jsem v~ukázkové implementaci použil zpětný log implementovaný zásobníkem.
Zpětný log vyžaduje při čtení kromě čtené hodnoty pouze kontrolu stavu příslušného
zámku, takže jde o~velmi rychlou operaci.

Použitá zásobníková
implementace je vhodná za předpokladu, že jedna transakce nebude opakovaně přepisovat
totéž místo. To je rozumný předpoklad, jelikož transakce by se celkově měly snažit
minimalizovat přístupy do sdílené paměti a~provádět výpočty v~rychlejší privátní
paměti.


\chapter{Ukládání na disk}

V~této kapitole nejprve prozkoumám úpravy, které je nutné provést v~implementaci
transakční paměti, aby mohl být průběh každé transakce zaznamenán, takže po restartu
aplikace může být průběh transakce \uv{přehrán}. Nejprve
popíši v~jakých formátech je možné ukládat obsah databáze na disk, a~poté jaký formát
jsem při implementaci zvolil.

\section{Rozšíření STM o~zápis změn}

Pro implementaci hybridní databáze je třeba STM rozšířit tak, aby při commitu transakce
zajistila zapsání změn na disk. Tedy doplnit logování v~průběhu transakce
o~informace, které umožní její opětovné provedení, a~upravit commit, aby tyto informace
zapsal na disk.

Zápis na disk může být pomalý a~vlákno s~transakcí nemusí chtít
čekat na jeho dokončení. Navíc zapisovaná data je nutno v~souboru seřadit (při použití
globálních hodin nejlépe dle času commitu). Proto je vhodné zavést pomocné vlákno (dále
servisní vlákno), které
bude od ostatních vláken dostávat data, která mají být zapsána, a~bude se starat
o~vlastní zápis. Data je nejlépe předávat pomocí globální fronty požadavků
(která se postará o~správné pořadí zápisu).

\subsection{Ukládání dat v~průběhu transakce}
Možností, jak držet potřebná data v~průběhu transakce, je několik. Nejpřímočařejší
je zapisovat do bufferu transakci tak, jak má být uložena
na disku. Výhodou je relativně snadná implementace (pokud je diskový formát vhodný),
naopak nevýhodnou může být prodloužení doby transakce (\uq{každý zápis se dělá dvakrát}).
Další variantou je místo zapisování logu v~nějakém vhodném interním formátu, který
budu převeden do formátu pro zápis až při / po (úspěšném) commitu. 

Toto řešení s~předchozím
sdílí problém vyššího množství zápisů do paměti, avšak při vhodném návrhu je výhodnější,
protože zkracuje čas strávený uvnitř transakce. Naopak zvětšuje čas pro zápis logu
na disk, ale to tolik nevadí, protože tato operace může probíhat již mimo transakční
paměť. Významnou výhodou obou těchto přístupů je nezávislost na implementaci
transakční paměti, takže mohou být použity i~v~případě, že implementace transakční
paměti je pevně daná -- což může být užitečné s~nástupem hardwarových implementací
transakční paměti.

Poslední variantou je integrace logu pro zápis na disk do standardního transakčního
logu, který je již součástí implementace transakční paměti. Výhodou tohoto přístupu
je snížení jak paměťové, tak časové režie, protože místo dvou logů obsahujících
z~velké části duplicitní informace (např. adresy, kam bylo zapisováno) máme log
pouze jeden. Podmínkou pro použití tohoto přístupu je samozřejmě možnost modifikovat
implementaci transakční paměti. Protože implementace transakční paměti je součástí
této práce, rozhodl jsem se pro tento přístup.

\subsection{Zpracování dat při commitu}
Pokud jsou potřebná data v~průběhu transakce ukládána do bufferu přímo v~diskovém
formátu, stačí je v~průběhu commitu zařadit do fronty ke zpracování servisním vláknem.
Pokud jsou ovšem uložena v~jiném tvaru, nabízí se otázka, zda je má do formátu
pro zápis převést vlákno vykonávající příslušnou transakci, nebo až servisní vlákno.
Otázka času zpracování se nabízí také pro uvolňování paměti (které vždy funguje
jako dopředný log).

Výhodou odsunutí práce do servisního vlákna je zrychlení commitu. To může vést k~vyššímu
výkonu aplikace, pokud zápisů není tolik, aby servisní vlákno zahltily (výrazného
zrychlení lze dosáhnout především ve chvíli, kdy zapisuje pouze jedno vlákno).
Druhou výhodou může být snazší implementace, protože paměť uvolňuje
pouze servisní vlákno (při vhodném řešení rollbacku), takže odpadají synchronizační
problémy.

Naopak výhodou ponechání práce na jednotlivých vláknech je vyšší propustnost servisního
vlákna. Problémem však může být komplikovanější synchronizace servisního vlákna,
jelikož frontu nestačí implementovat jako spojový seznam (chráněný zámkem při přidávání
položek vlákny transakcí) se semaforem, který určuje počet nezpracovaných položek v~něm,
ale je nutné nejprve transakci zařadit do fronty, poté zpracovat transakční log
a~nakonec oznámit servisnímu vláknu, že transakce je připravena k~zápisu.
(Složitější synchronizaci se lze vyhnout zpracováním logu pod ochranou zámku fronty,
což je ale nevhodné, protože to výrazně prodlouží dobu držení zámku, nebo zpracováním
na počátku commitu před přidáním do fronty, což zvyšuje riziko
kolize s~jinou transakcí.)


%%%%%%%%%%%%%%%%%%%%%%%%%%%%%%%%%%%%%%%%%%%%%%%%%%%%%%%%%%%%%%%%%%%%%%%%%%%%%%%%

\section{Zvolený diskový formát}

V~této sekci popíši možnosti, jak ukládat databázi na disk, a~jaké řešení jsem
zvolil pro hybridní databázi.

\subsection{Datové formáty obecně}
Většina databází (především relačních) používá souborový formát respektující strukturu dat
-- typicky soubory odpovídají tabulkám a~každý soubor je vnitřně rozdělen na části
stejné velikosti, které odpovídají jednotlivým řádkům (v~praxi používané formáty
jsou samozřejmě složitější kvůli podpoře pro data neomezené délky, přímo do datového
souboru zakomponované indexy atp.). Výhodou je jednoduchá práce s~daty, protože
ta lze přímo načítat a~zapisovat do souboru bez složitého zpracování v~paměti. Naopak
problémem je nízká rychlost zápisu, protože změny v~různých řádcích databáze znamenají
zápisy na různá místa různých souborů, což je u~rotačních disků velmi pomalé a~u~ssd
disků může velmi snižovat životnost (protože disky SSD jsou schopny přepisovat data
pouze po velkých blocích (řádově desítky až stovky kilobytů \cite{ssd})
a~každý blok vydrží jen omezené množství přepsání a~i~sebemenší změna v~databázi přepíše
celý blok).

Ještě významnějším problém je nemožnost zaručení konzistence dat, protože
zápis změn zjevně není atomickou operací. Proto se většinou tento způsob uložení dat
doplňuje ještě o~transakční log, kam se zapisují operace, které se na databázi provádějí.
S~jeho pomocí je v~případě havárie možné převést data zpět do posledního uloženého
konzistentního stavu. Užití logu může mít i~pozitivní dopad na výkon. Celková propustnost
se sice sníží, protože je kromě změn přímo v~datových souborech je navíc zapisován
i~log, ale transakci je při vhodné implementaci možné považovat za úspěšnou již
po jejím zapsání logu, které je rychlejší než modifikace datových souborů, protože
do logu se zapisuje sekvenčně, takže latence při commitu transakce může výrazně klesnout.

Datový formát založený na logu změn je naopak oblíben u~noSQL in-memory databází
(např. Redis \cite{Redis}).
Jeho výhodou je snadný zápis a~rychlý zápis. Naopak problémem může být nutnost
zpracovat celý log dříve, než můžeme číst data (což příliš nevadí, čteme-li
log pouze při startu databáze). Závažnější je zvětšování logu v~průběhu práce
s~databází, protože do logu se data pouze připisují a~nikdy se neodebírají.
Běžným řešením je při překročení určité velikosti logu vygenerovat log
nový, který obsahuje nejkratší posloupnost operací vedoucí k~současnému stavu databáze,
a~jím nahradit ten starý, příliš dlouhý. Aby tento přístup fungoval, nesmí
být množství zápisů do databáze příliš velké.


\subsection{Formát Hybridní databáze}
Protože Hybridní databáze uchovává za běhu všechna data v~paměti, zvolil jsem
diskový formát založený na logu.
Navržený diskový formát se skládá ze dvou vrstev. Spodní vrstva se stará o~detaily
uložení na disku a~poskytuje vyšší vrstvě jednoduchou abstrakci pro strukturovaná data.
Vyšší vrstva pak definuje prostředky té nižší, jak jsou uloženy jednotlivé uzly
a~operace s~nimi.

\subsubsection*{Nižší vrstva}

Cílem spodní vrstvy je poskytnout vyšší vrstvě abstrakci pro strukturovaná
data, v~jejímž rámci je pak vybudován vlastní formát databáze, při zachování maximální
jednoduchosti a~tím i~rychlosti zpracování. Každý datový soubor se skládá z~jednotlivých
tzv. kořenových elementů. Kořenovým elementem může být pole (seznam) či řetězec.
Řetězcem je zde myšlena konečná posloupnost bytů (tedy speciálně může obsahovat znak
\verb|'\0'|) a~pole je konečná posloupnost jiných polí a~řetězců. 

Klíčovou odlišností
kořenových elementů oproti elementům vnořeným je, že před každým z~nich je uvedena
celková délka daného elementu a~na konci každého z~nich může být CRC kód.
Pro přehlednost uveďme gramatiku popisující nižší vrstvu, přičemž
přesný způsob uložení jednotlivých částí bude popsán dále:

\begin{alltt}
{\sl{}soubor} := {\sl{}kořenový\_element}*
{\sl{}kořenový\_element} := '[' "délka kořenového elementu" ']'
                    {\sl{}element} ( '#' "CRC32 Castagnoli" )? '.'

{\sl{}element} := {\sl{}řetězec} | {\sl{}pole}
{\sl{}pole}    := '(' {\sl{}element}* ')'
{\sl{}řetězec} := "délka řetězce" "obsah řetězce"
\end{alltt}


Délka kořenových elementů je užitečná spíše pro snadné načítání, kdy je celý kořenový
element načten do bufferu v~paměti, takže při čtení jednotlivých řetězců může být
předáván přímo ukazatel do tohoto bufferu. CRC (Cyclic Redundacy Check) je druh
kontrolního součtu, který umožňuje odhalovat chyby při přenosu
či ukládání dat. Jeho přítomnost v~kořenových elementech je důležitá pro
implementaci transakcí. Každá transakce je uložena jako pole na kořenové úrovni,
takže CRC kód na jejím konci lze použít k~ověření, že byla korektně uložena.
Z~různých verzí CRC byla zvolena verze CRC32 Castagnoli \cite{CRC32c},
která je použita
kupř. v~protokolu iSCSI \cite{iSCSI} či souborovém systému Btrfs \cite{btrfs}.
Navíc je implementace této varianty CRC jednou z~instrukcí
SSE4.2 \cite{sse42},
takže v~moderních procesorech může být výpočet CRC kódu výrazně urychlen.
(Ač tato optimalizace není v~současné implementaci obsažena.)
Měření totiž ukazují, že výpočet CRC je, po práci s~diskem samozřejmě, časově
nejnáročnější částí práce s~daty v~tomto formátu.

Co se týče způsobu uložení jednotlivých prvků, tak délka kořenového elementu a~CRC kód
jsou vždy uloženy jako 8-bajtové resp. 4-bajtové celé číslo v~little-endian tvaru.
Zajímavější je uložení řídicích znaků \verb|[]()#.| a~délek řetězců. Řídicí znaky
jsou zobrazeny na celá čísla (tzv. řídicí kódy) z~rozsahu 0 až 15 (interval je větší
než počet řídících
znaků pro případ, že by se ukázalo jako vhodné formát nějak rozšířit) a~délky řetězců
na řídicí kódy od 16 výše (tedy např. řetězec délky 10 bude mít kód 26). Díky tomuto 
přístupu se délky řetězců nijak neodlišují od jiných řídicích znaků a~zjistit další
element lze načtením jednoho řídicího kódu. 

Rozsah řídicích kódů je 0 až $2^{64} - 1$.
Pro úsporu místa však nejsou ukládány s~pevnou šířkou, ale způsobem připomínajícím
UTF-8 \cite{utf8}, např. kódy menší než 128 jsou uloženy v~jednom bajtu způsobem
\verb|0b0xxxxxxx|, kódy od 128 do $2^{14} - 1$ ve tvaru \verb|0b10xxxxxx xxxxxxxx| atd.
Tedy počet jedniček na začátku prvního bajtu je o~jedna menší, než kolik bajtů celý
řídící kód zabírá. Pokud je těchto jedniček méně než osm, následuje po nich nula
a~až za ní vlastní kód. Je-li první bajt \verb|OxFF|, tak žádná nula vložena není
a~následuje přímo osm bajtů kódu. Do vyhrazených bitů je kód uložen následovně:
Nejvyšší bity jsou uloženy v~prvním bajtu a~zbylé jsou uloženy do následujících bajtů
little-endian způsobem.

\subsubsection*{Vyšší vrstva}
Vyšší vrstva obsahuje popis formátů databáze prostřednictvím prostředků nižší vrstvy.
Databáze se skládá z~několika souboru, přičemž jméno každého z~nich se skládá ze
jména konkrétní databáze a~přípony. Vždy přítomný je soubor s~příponou \verb|.schema|,
který obsahuje popis databáze, tj. název typu databáze, seznam typů uzlů a~u~každého
typu uzlu popis jeho atributů. Při otevírání databáze je tento soubor porovnán
s~popisem příslušné databáze v~paměti, aby se předešlo otevření nekompatibilní databáze.

Ostatní soubory jsou datové a~mají jako příponu přirozené číslo. Absolutní hodnota
čísla není důležitá, slouží pouze k~seřazení souborů do správného pořadí a~zjištění,
zda v~posloupnosti souborů nejsou díry. Přípona není na detekci správného
pořadí dostatečná, jelikož je snadné soubory pomíchat, především ve~chvíli, kdy se
je z~snažíme přejmenovat tak, aby začínaly od jedné. Proto obsahuje každý datový soubor
na začátku a~na konci magický řetězec. Magický řetězec je osm bajtů dlouhý a~náhodně
generovaný, přičemž jeho hodnota se na konci jednoho souboru a~začátku na něj
navazujícího souboru shoduje. Tím je zajištěno (s~rozumnou pravděpodobností), že se
soubory nepomíchají.

Uložení dat ve více souborech je nutné především kvůli možnosti zkrátit log
vypsáním databáze do nového souboru. Výhodou
takového uspořádání je také možnost velmi snadného vytváření snímků stavu databáze
v~určitém čase -- ve chvíli, kdy chceme vytvořit snímek, založíme nový datový soubor
a~všechny následující transakce budou zapisovány do něj.

Samotný formát datových souborů je velmi prostý. Prvním kořenovým elementem je pole
obsahující typ databáze, k~níž soubor patří, pak následuje řetězec obsahující magický
řetězec a~případně řetězec označující začátek výpisu databáze. První načítaný soubor
vždy musí obsahovat výpis databáze. Poté následují jednotlivé
transakce, přičemž každá transakce je uložena jako pole na kořenové úrovni. Prvky
transakce jsou opět pole. Každé z~nich popisuje jednu operaci s~databází a~to
následujícím způsobem. Prvním prvkem je řetězec popisující druh operace (možné operace
jsou vytvoření či smazání uzlu a~změna hodnoty atributu uzlu), za ním následuje
identifikátor uzlu (celé kladné číslo). Dále podle druhu operace buď číslo typu uzlu
(pro vytvoření nového uzlu), nic (při mazání), nebo pořadí upravovaného atributu
a~jeho nová hodnota. Pokud je soubor ukončený (což jsou všechny kromě posledního),
obsahuje na konci ukončovací řetězec a~magický řetězec navazujícího souboru.

Formát uložení jednotlivých typů atributů je následující: 
Řetězce (\verb|String| a~\verb|RawString|) jsou uloženy jako řetězce stejně jako v~paměti,
pouze ukončovací \verb|'\0'| u~typu \verb|String| se do souboru neukládá.
Znaménkové číselné typy
(\verb|Int{8,16,32,64}|) a~čísla s~plovoucí řádovou čárkou (\verb|Float|,
\verb|Double| a~\verb|LDouble|) jsou uložena jako řetězce tak,
jak je architektura AMD64 ukládá
v~paměti (tedy little-endian tvar a~čísla s~plovoucí řádovou čárkou
jsou v~tvaru dle normy \hbox{IEEE 754}
\cite{IEEE754}). 
Bezznaménková celá čísla také zachovávají little-endian tvar, ale pro úsporu místa
se nejvyšší nulové bajty neukládají. Ukazatele na uzly jsou na disku reprezentovány
identifikátorem uzlu, na nějž odkazují. Ten je uloženo jako bezznaménkové celé číslo.

\subsection{Vypsání databáze}
Důležitou součástí každého formátu databáze založeného na logu je způsob, jakým
řeší nahrazení příliš dlouhého logu novým, krátkým se stejným obsahem. První možností
by bylo ve chvíli, kdy se rozhodneme, že současný log je příliš dlouhý, zastavit
databázi, vypsat její obsah do nového souboru a~pokračovat v~práci. Problém toho
přístupu je, že čas strávený výpisem databáze může být dlouhý a~nemusí být žádoucí
aplikaci na tuto dobu zastavit.

Druhou možností, použitou např. v~databázi Redis \cite{Redis}, je provést systémové
volání \verb|fork|.
To \uv{rozdělí} současné vlákno na dvě, přičemž nově vytvořená kopie se stane
samostatným procesem. Tento proces má zkopírované všechny zdroje (paměť
i~otevřené file deskriptory) svého rodiče. Může tedy vypsat databázi, protože
vlastní snímek paměti rodiče ze chvíle, kdy bylo volání \verb|fork| vykonáno.
Poté pouze informuje rodiče, že vypsal databázi, ten do ní dopíše kousek logu,
co vznik v~čase, kdy výpis probíhal, a~pokračuje v~práci s~novým logem.
Nevýhody \verb|fork|u spočívají právě v~jeho schopnosti vyrobit snímek zdrojů
rodičovského procesu. Kromě potenciálně velkého počtu zkopírovaných file deskriptorů,
které musí potomek uklidit, ač je vůbec nepotřeboval, může nastat problém s~nedostatkem
paměti. 

Linux implementuje \verb|fork|, pomocí přístupu copy-on-write
\cite{linuxCOW}, takže samotné volání je velmi rychlé a~efektivní. Pokud ale rodič
začne do paměti více zapisovat, jsou měněné stránky kopírovány, protože jsou
sdíleny s~potomkem, který v~nich pořád musí vidět původní data. V~extrémním případě
tedy potomek spotřebovává tolik paměti, co jeho rodič. Že by se toto stalo s~pamětí
užívanou hybridní databází je nepravděpodobné, ale velmi snadno může tyto problémy
způsobit část programu, která s~databází vůbec nesouvisí a~pouze hodně zapisuje do paměti.
V~případě
Redisu to nevadí, protože běží jako samostatný proces a~má tedy svou paměť 
i~deskriptory pod kontrolou.
Naopak databáze implementovaná jako knihovna funkcí nemá žádnou kontrolu nad prostředky,
které alokuje zbytek aplikace, takže využití \verb|fork|u je nevhodné.

Proto jsem zvolil komplikovanější přístup, který se ale obejde jak bez zastavení
práce s~databází, tak bez volání \verb|fork|. Když chceme vypsat databázi,
tak je založen nový datový soubor a~na jeho začátek je zapsáno, že obsahuje výpis
databáze. Nové transakce se zapisují do tohoto nového souboru. Výpis databáze
vypadá jako speciální druh transakcí (v~souborovém formátu má jako první prvek
v~poli řetězec nulové délky), které jsou zamíchány mezi běžné transakce. Tyto speciální
transakce obsahují výpisy jednotlivých uzlů (současná implementace obsahuje v~každé
transakci výpis jednoho uzlu). Výpis uzlu je pole, v~němž je identifikátor typu uzlu
(jeho pořadí v~rámci typu databáze, číslováno od 0), 
identifikátor uzlu
a~za nimi následují hodnoty atributů. Když jsou vypsány všechny uzlu, je do souboru
zapsána značka \uv{konec výpisu}, takže při načítání databáze lze určit, zda vypsání
proběhlo úplně, nebo bylo přerušeno.

Samotné vypisování uzlů je prováděno servisním vláknem. To ve chvíli, kdy by mělo
z~fronty vybrat další požadavek, zkontroluje, zda neprobíhá výpis databáze.
Pokud probíhá,
tak provede vypsání přednastaveného počtu uzlů (v~základním nastavení 5,
ale počet je možné změnit
konfigurační volbou), a~až pak provede další událost čekající ve frontě. Je-li fronta
prázdná, vypíše další uzly a~opět zkontroluje frontu. Drobným zádrhelem, na který je
dobré pamatovat, je, že současná implementace při výpisu obchází frontu transakcí
(což je z~hlediska výkonu dobře),
takže výpis uzlu se v~souboru může objevit před transakcí, která ho vytvořila.
Proto kód načítající databázi musí být na tuto možnost připraven a~správně ji ošetřit,
tedy požadavek na vytvoření existujícího uzlu ignorovat a~ne vyvolat chybu.


\chapter{Rozhraní databáze}
Jak bylo zmíněno výše, hybridní databáze je implementována jako knihovna
funkcí v~jazyce C. Protože jazyk C neumožňuje provést definice typů uzlů,
indexů a~databází ve tvaru rozumně zpracovatelném pro člověka i~počítač,
jsou tyto definice provedeny ve speciálním, k~tomuto účelu vytvořeném, jazyce.
Tyto definice jsou následně přeloženy do jazyka C. (Při implementaci překladače
jsem použil
jazyk Perl \cite{Perl} s~knihovnou RecDescent \cite{RecDescent}.)
Popis návrhu hybridní databáze začneme popisem rozhraní z~hlediska běžného uživatele,
poté představíme jazyk pro generování definic. V~příloze 1 naleznete jednoduchý
příklad použití databáze.

\section{Hybridní databáze z~pohledu uživatele}

Nejprve popíšeme rozhraní, které bude uživatel používat při běžné práci s~databází.
Datové typy \verb|Node|, \verb|Handle| a~\verb|Database| jsou abstraktní
třídy. Proto všechny funkce, které je mají jako parametr,
jsou v~implementaci obaleny makry, takže bez explicitního
přetypování je jim jako parametry možné předat i~podtřídy daných tříd.
Pokud je nějaká funkce má jako návratovou hodnotu, většinou existuje makro
s~jiným jménem obalující funkci tak, aby přímo vracela správný typ (často
se parametry makra mírně liší od funkce, protože kupř. makra berou jako parametry
místo deskriptorů typů jejich názvy).

\subsection{Vytváření a~zánik databáze}
\begin{verbatim}
Database *database_create(const DatabaseType *type,
                          const char *file,
                          enum DbFlags flags);
#define dbCreate(Type, file, flags)

void database_close(Database *D);
\end{verbatim}

Databázi je možné vytvořit voláním funkce \verb|database_create| nebo makra
\verb|dbCreate|. Výhodou makra je, že jeho návratový typ není \verb|Database*|,
ale ukazatel na konkrétní podtřídu typu \verb|Database| (uvnitř makro stejně volá funkci
\verb|database_create|).

Funkce \verb|database_create| nejprve alokuje a~inicializuje objekt databáze.
Poté zkontroluje schéma databáze \verb|file.schema| (tj. soubor obsahující popis typů
uzlů a~jejich atributů, které databáze může obsahovat) a~načte data uložená v~souborech
\verb|file.N|, kde \verb|N| je přirozené číslo. Čísla souborů nemusí začínat od
jedné, ale posloupnost nesmí obsahovat mezeru, jinak vytváření databáze skončí chybou.
(Protože dává smysl smazat první soubor za předpokladu, že databáze byla vypsána
do druhého souboru, ale není možné odstranit nějaké transakce uprostřed historie.)
První soubor v~posloupnosti také musí obsahovat výpis databáze.
(Podobnější informace o~diskovém formátu naleznete ve třetí kapitole.)
Když jsou všechna data v~paměti, proběhne inicializace indexů a~na závěr jsou spuštěna
pomocná vlákna.

Pomocná vlákna jsou dvě -- servisní vlákno a~vlákno časovače. Servisní vlákno se stará
o~zápis transakcí na disk. Ty jsou mu předávány pomocí globální fronty jednotlivými
transakčními vlákny. (Míra zpracování transakcí v~okamžiku jejich předání servisnímu
vláknu závisí na konfiguračních volbách databáze.) Kromě toho ještě zajišťuje
vypsání databáze na disk a~zakládání nových souborů. Servisní vlákno transakce
nezapisuje okamžitě (není-li to explicitně vyžadováno), ale ukládá je do bufferu,
který je zapsán, až když je plný. Aby bylo zajištěno, že i~pří nízkém množství zápisů
budou data zapsána v~konečném čase, existuje vlákno časovače. Vlákno časovače periodicky
(v~základním nastavení každých 5~s) posílá zprávu servisnímu vláknu, aby zapsalo obsah
bufferu na disk a~zavolalo funkci \verb|fsync|. Ta přinutí systém data skutečně
zapsat, a~ne držet je ve vlastních bufferech.

Databáze se uzavírá voláním \verb|database_close|. Je na
uživateli databáze, aby zajistil, že v~době tohoto volání již žádné vlákno s~databází
nepracuje.

\subsection{Handly a~transakce}
\begin{verbatim}
#define dbHandleCreate(D)
Handle *db_handle_create(Database *D);

void db_handle_free(Handle *H);

void tr_begin(Handle *H);
void tr_abort(Handle *H);
bool tr_commit(Handle *H, enum CommitType commit_type);

void tr_hard_abort(Hadler *H);
\end{verbatim}

Dříve než může vlákno začít transakci, potřebuje handle. Handle představuje
jakýsi kontext transakce -- obsahuje transakční log, readset a~všechny další
nezbytnosti pro její běh. Transakce samy nemají podobu objektů, ale vždy
je funkcím v~rámci transakce předáván handle. Výhodou je, že jeden handle může
být postupně využit pro mnoho transakcí, čímž se šetří opakované alokace paměti,
a~protože databáze potřebuje mít přehled o~běžících transakcích, i~množství
práce s~centrální strukturou.

Funkce \verb|tr_begin| zahájí novou transakci. Co skutečně vykoná závisí na tom,
zahajujeme-li hlavní transakci nebo běží-li již nějaká transakce a~zahajujeme pouze
vnořenou transakci. Práce s~hlavní transakcí závisí na implementaci STM,
která byla popsána ve druhé kapitole, takže se ji zde nebudu
věnovat. Naopak vnořené transakce se odehrávají zcela v~rámci příslušné hlavní
transakce a~fungují jako \uv{záložky} v~transakčním logu -- když je zahájena vnořená
transakce, je na zásobník vnořených transakcí, který každý handle obsahuje, umístěna
položka popisující současný stav transakčního logu a~readsetu.

Obdobně při úspěšném ukončení transakce pomocí \verb|tr_commit| nebo jejím zrušení
voláním \verb|tr_abort|. Při ukončení vnořené
transakce pomocí \verb|tr_commit|, volání \verb|tr_commit| vždy uspěje a~pouze
odstraní položku z~vrcholu zásobníku vnořených transakcí, čímž se daná vnořená transakce
stane součástí nadřazené transakce. Parametr \verb|commit_type| určuje, zda bude commit
synchronní, tj. funkce \verb|tr_commit| se vrátí až po zapsání celé transakce na disk,
či asynchronní, kdy návrat z~funkce proběhne okamžitě po předání dat servisnímu
vláknu.

Funkce \verb|tr_abort| a~\verb|tr_hard_abort| provedou rollback transakce (tedy vrátí
změny, které transakce doposud provedla). Rozdíl mezi nimi je ten, že \verb|tr_abort|
zruší pouze aktuální (vnořenou) transakci, zatímco \verb|tr_hard_abort| vždy zruší celou
hlavní transakci. To je užitečné pokud nějaká operace uvnitř transakce selže z~důvodu
kolize s~jinou transakcí. V~ideálním světě by v~tu chvíli příslušná operace sama
provedla rollback a~restart transakce (například vyvoláním k~tomu určené výjimky).
Jelikož ale jazyk C neobsahuje výjimky a~použití nějaké konkrétní implementace výjimek
by uživatele nutilo používat danou implementaci v~celém jeho kódu, vracejí všechny
funkce, které mohou selhat, hodnotu typu \verb|bool| přičemž \verb|true| značí úspěch
a~\verb|false| selhání. Při selhání je pak na uživateli, aby zajistil zrušení transakce
voláním \verb|tr_hard_abort| a~její následné restartování.

\subsection{Makra pro usnadnění práce s~transakcemi}
\begin{verbatim}
#define trFail

#define trBegin
#define trCommit(commit_type, restart_action)
#define trAbort
\end{verbatim}

Pro pohodlnější práci s~transakcemi jsme implementovali i~makra, která doplňují jistou
podporu pro automatické restartování transakcí. Makra \verb|trBegin| a~\verb|trCommit|
spolu tvoří závorky obalující transakci. Pokud selže commit nebo nějaká operace uvnitř
transakce, tak se transakce sama zruší a~restartuje. Aby se předešlo live-locku,
transakce před restartem počká náhodný čas, přičemž maximální čas čekání exponenciálně
roste v~závislosti na tom, o~kolikátý restart dané transakce se jedná.

Jak tato makra zjistí, že operace uvnitř transakce selhala, když nemají k~dispozici
výjimky? Jednoduše -- každé makro při selhání vykoná makro \verb|trFail|. To je
defaultně nastaveno jako skok na návěstí \verb|tr_failed|, které je definované uvnitř
makra \verb|trCommit|. Tento přístup má samozřejmě svá omezení, např. nefunguje
transparentně, pokud práce s~databází probíhá uvnitř volání funkce. Naopak ale umožňuje
uživateli přizpůsobit ho svým potřebám předefinováním makra \verb|trFail|. Z~důvodu
jisté omezenosti toho přístupu má makro \verb|trCommit| jako druhý, volitelný parametr
kousek kódu, který je proveden, pokud daná transakce selže a~není transakcí hlavní.
V~tom případě nemůže být přímo restartována, ale musí propagovat selhání až do hlavní
transakce.

Makra \verb|trBegin|, \verb|trCommit| a~\verb|trAbort| pracují s~handlem jménem
\verb|H|. Pokud je nutné použít jiný handle, existují k~těmto makrům varianty,
které mají na konci jména podtržítko a~mají navíc jako první parametr handle.

\subsection{Práce s~uzly a~indexy}
\begin{verbatim}
#define trNodeCreate(NodeType)
#define trNodeDelete(node)

#define trRead(node, Attribute)
#define trWrite(node, Attribute, new_value)

#define trUpdateIndexes(node)

#define nodeCast(NodeType, node)
\end{verbatim}

Stejně jako výše uvedená makra, i~tato pracují s~handleem jménem \verb|H| (kromě makra
\verb|nodeCast|) a~při selhání provedou makro \verb|trFail|. Makro \verb|trNodeCreate|
vytvoří nový uzel daného typu a~aktualizuje příslušné indexy. Obdobně makro 
\verb|trNodeDelete|
uzel smaže a~provede aktualizaci indexů. Naopak makro \verb|trWrite|, které slouží pro
zápis do existujícího uzlu, po svém provedení indexy neaktualizuje.

K~tomuto chování
má dva důvody: Za prvé není neobvyklé, že v~rámci jedné transakce je jednomu uzlu
modifikováno několik atributů, a~pokud by indexy byly aktualizovány po zápisu každého
z~nich, představovalo by to zbytečnou režii. Druhý důvod je možná ještě důležitější
-- pro indexy může být výhodné předpokládat, že každý uzel splňuje určité invarianty,
čehož není možné dosáhnout, je-li index aktualizován po každém zápisu, protože přechod
mezi korektními stavy může vyžadovat modifikaci několika atributů. Proto existuje
makro \verb|trUpdateIndexes|, které informuje indexy o~změně příslušného uzlu.
Toto makro je nutné zavolat na každý modifikovaný uzel před koncem příslušné transakce.

Makro \verb|trRead| slouží ke čtení hodnot atributů uzlů. Jeho důležitou vlastností
je, že poskytuje konzistentní pohled na data (problematika konzistentního čtení
je důkladněji probrána v~následující kapitole), takže uživatel je odstíněn od změn
prováděných jinými transakcemi. 

\subsection{Práce s~indexy}

\begin{verbatim}
#define trIndex(jméno_indexu, metoda_index, parametry...)
\end{verbatim}
Jak bylo zmíněno výše, indexy jsou klíčem pro přístup k~jednotlivým uzlům.
V~rámci této práce rozumíme indexem libovolnou datovou strukturu, jejíž
stav je jednoznačně určen stavem databáze, nad níž je vytvořena. Protože jsou indexy,
jako jediný způsob pro (počáteční) přístup k~uzlům databáze, pro práci s~ní
klíčové, byly navrhovány se snahou vytvořit je co nejobecnější při zachování
jednoduchosti práce.

Z~hlediska databáze
je rozhraní indexů velmi jednoduché -- každý index se skládá z~kontextu, tedy vlastní
datové struktury (která je databází automaticky alokována) a~callback funkce, která
je volána při změně (relevantní části) databáze. Callback má jako parametry kontext,
typ události, uzel, k~němuž se daná událost váže, a~samozřejmě handle, protože update
indexů vždy probíhá v~rámci příslušné transakce. Protože je i~celý kontext indexu
součásti transakční paměti, chová se z~pohledu ostatních vláken stejně jako zbytek
databáze -- především tedy ostatní vlákna vždy vidí index v~konzistentním stavu.
Kromě callback funkce obsahuje každý index ještě funkce pro inicializaci kontextu
a~uvolnění prostředků indexem užívaných.

Z~pohledu uživatele je rozhraní indexu zcela v~režii jeho tvůrce. K~jednotlivým metodám
se přistupuje pomocí makra \verb|trIndex|. Toto makro obaluje skutečnou funkci, která 
má hlavičku:
\begin{verbatim}
bool jméno_metody(const Typ_kontextu* context, Handle *H,
                  NávratovýTyp *value, další_parametry...)
\end{verbatim}

Návratová hodnota je vždy typu \verb|bool| a~signalizuje, zda bylo volání úspěšné,
nebo došlo ke kolizi a~je nutné transakci restartovat (\verb|true| značí úspěch).
Toto chování je uživateli neviditelné. V~případě selhání makro \verb|trIndex| provede
\verb|trFail| a~bude tak selhání propagovat dále. V~případě úspěchu vrátí hodnotu
\verb|*value|. Z~pohledu tvůrce indexu stojí za pozornost, že uvnitř definice metody
indexu je makro \verb|trFail| definováno jako \verb|return false| místo běžné definice
\verb|goto tr_failed|.


\section{Generování definic}
Součástí ukázkové implementace je i~jazyk pro deklarování typů indexů, uzlů a~databází.
Původně jsem se chtěl obejít bez tvorby nového jazyka, ale ukázalo se, že tyto
typy není možné dost přehledně deklarovat pouze pomocí jazyka C a~jeho vestavěného
preprocesoru. Proto jsem vytvořil jazyk pro tyto definice. Zdrojový kód v~tomto jazyce
je překládán pomocí skriptu v~jazyce Perl \cite{Perl} do jazyka C.
Kromě deklarací typů užívám obdobný jazyk i~pro generování funkcí pro práci
s~jednotlivými typy atributů.
Perlový překladač staví na knihovně RecDescent \cite{RecDescent}, která umožňuje
snadné parsování textu na základě jeho gramatiky. 
Nejprve ukážeme základní syntaxi navrženého jazyka, a~poté
popíšeme jednotlivé typy definic.

\subsection{Základní syntaxe}
\begin{alltt}
{\sl{}file} := ( {\sl{}interface} | {\sl{}implementation} | {\sl{}include} | {\sl{}node\_type} | 
        {\sl{}index\_type} | {\sl{}database\_type} )*

{\sl{}interface} := 'Interface' {\sl{}block}
{\sl{}implementation} := 'Implementation' {\sl{}block}
{\sl{}include} := 'Include' /"([^"]+)"/
\end{alltt}

Nejprve pár slov k~syntaxi. Jazyk rozlišuje velikost písmen.
Parser zpracovává kód po tokenech, které mohou být odděleny
libovolným počtem bílých znaků. Za bílý znak je považován i~středník a~komentáře,
které začínají znakem~\verb|#| a~pokračují do konce řádku.
Element \verb|block| představuje blok kódu uzavřený ve složených závorkách. Ty se
mohou uvnitř bloku vyskytovat za předpokladu, že jsou spárované. Celý blok je považován
za jeden token.

Element \verb|Include| slouží k~načtení obsahu jiného souboru. Jméno souboru v~uvozovkách
je interpretováno jako relativní či absolutní, nedochází k~prohledávání žádných
adresářů s~hlavičkovými soubory. Při vkládání souboru není kontrolováno opakované vložení
souboru. Sémanticky odpovídá direktivě \verb|#include|
jazyka C. Elementy \verb|Interface| a~\verb|Implementation| slouží k~vložení bloku
kódu v~jazyce C před vygenerované deklarace resp. definice. (Tedy pokud v~implementaci
indexu budu chtít použít funkce z~hlavičkového souboru \verb|foo.h|, tak do kódu
vložím \verb|Implementation { #include "foo.h" }|.) Těchto elementů se v~jednom souboru
může nacházet libovolné množství, přičemž ve výsledném souboru jsou za sebou tak,
jak byly uvedeny.

\subsection{Definice uzlů}
\begin{alltt}
{\sl{}node\_type} := 'NodeType' {\sl{}name} '\{' {\sl{}node\_attribute}* '\}'
{\sl{}node\_attribute} := 'Attribute' {\sl{}name} ':' {\sl{}attribute\_type}
\end{alltt}

Typy uzlů se definují pomocí klíčového slova \verb|NodeType|. Definice pak obsahuje
jména atributů a~jejich typy, přičemž na pořadí atributů záleží.
Současně implementované jsou následující typy atributů
(které pokrývají základní typy jazyka C):

\begin{itemize}
  \item Znaménkové typy \verb|Int8|, \verb|Int16|, \verb|Int32| a~\verb|Int64|
    jsou v~paměti realizovány datovými typy \verb|int{8,16,32,64}_t|.
  \item Neznaménkové typy \verb|UInt8|, \verb|UInt16|, \verb|UInt32| a~\verb|UInt64|
    jsou v~paměti realizovány datovými typy \verb|uint{8,16,32,64}_t|.
  \item Reálná čísla \verb|Float|, \verb|Double| a~\verb|LDouble| jsou realizovány
    datovými typy \verb|float|, \verb|double| a~\verb|long double|.
  \item Typ \verb|String| představuje textový řetězec ukončený znakem \verb|'\0'|
    a~je realizován datovým typem \verb|char*|.
  \item Typ \verb|RawString| představuje řetězec, který může obsahovat hodnotu
    \verb|'\0'|. V~paměti je reprezentován typem \verb|struct RawString|.
  \item Typ \verb|Pointer| představuje ukazatel na jiný uzel. Z~technických důvodů
    je realizován jako \verb|union| identifikátoru uzlu a~ukazatele na uzel. Z~pohledu
    uživatele se jeví jako pouze jako \verb|Node*|. Použití \verb|union|u je nutné,
    protože při načítání dat ze souboru je do atributu typu \verb|Pointer|
    nejprve uložen identifikátor uzlu, na který ukazuje, a~až když jsou načteny
    všechny uzly,
    je možné je projít a~změnit identifikátory uzlů na ukazatele na ně.
\end{itemize}


\subsection{Definice indexů}
\begin{alltt}
{\sl{}index\_type} := 'IndexType' {\sl{}name} ':' {\sl{}type\_of\_context}
              '\{' {\sl{}i\_init i\_destroy i\_update i\_method}* '\}'
{\sl{}i\_init} := 'Init' {\sl{}block}
{\sl{}i\_destroy} := 'Destroy' {\sl{}block}
{\sl{}i\_update} := 'Update' {\sl{}block}
{\sl{}i\_method} := 'Method' {\sl{}name} ( '(' {\sl{}arguments...} ')' )? ':' 
            {\sl{}return\_type block}
\end{alltt}

Každý typ indexu se skládá alespoň z~typu kontextu, konstruktoru, destruktoru
a~metody \verb|Update|. Dále by měl obsahovat alespoň jednu metodu, protože
metody indexu jsou jediným způsobem, jak z~něj získat informace.
Typ kontextu je datový typ, který bude sloužit jako kontext 
pro příslušný typ indexu,
tj. ukazatel na objekt toho typu bude předáván všem metodám indexu.

Konstruktor a~destruktor (tedy metody \verb|Init| a~\verb|Destroy|) jako jediné
pracují mimo transakce, protože pracují ještě před startem, resp. po zastavení
transakční paměti. Obě mají jako parametr kontext a~\verb|GenericAllocator|, pomocí
něhož mohou alokovat, resp. uvolňovat paměť. K~práci nemohou použít \verb|malloc|
a~\verb|free|, protože celý index musí být součástí transakční paměti, aby
mohl správně fungovat případný rollback transakce.

Metoda \verb|Update| slouží k~aktualizaci indexu a~je kromě konstruktoru a~destruktoru
jedinou metodou, která smí upravovat kontext. Jejími parametry jsou kontext, handle
transakce, druh události a~uzel, jehož se událost týká. Událost smazání uzlu je
provedena dříve, než je příslušný uzel smazán, takže metoda \verb|Update| s~ním
může pracovat. Naopak při vytvoření uzlu je událost vyvolána, až když je uzel vytvořen
a~inicializován. V~metodě \verb|Update| lze k~práci s~pamětí používat funkce
\verb|tr_memory_alloc()| a~\verb|tr_memory_free()|. Obě mají jako první parametr
handle.

Uživatelské metody vytvořené pomocí klíčového slova \verb|Method| mají jako parametry
kontext, handle transakce, ukazatel na návratovou hodnotu a~uživatelem definované
argumenty. Uživatelské metody stejně jako \verb|Update| vždy probíhají v~rámci transakce
a~mají návratovou hodnotu typu \verb|bool|, která značí, zda metoda byla úspěšná
(\verb|true|), nebo došlo ke kolizi a~je třeba restartovat transakci (\verb|false|).
Proto uživatelské metody nemohou vracet žádnou další informaci přímo, ale je nutné
využít parametru typu ukazatel.

\subsection{Definice databází}
\begin{alltt}
{\sl{}database\_type} := 'DatabaseType' {\sl{}name} '(' {\sl{}version} ')' 
                 '\{' ( {\sl{}node} | {\sl{}index} )? '\}'
{\sl{}node} := 'NodeType' {\sl{}name}
{\sl{}index} := 'Index' {\sl{}name} 'for' ( {\sl{}node\_type} | 'all' ) ':' {\sl{}index\_type}
\end{alltt}

Databáze je určena jménem, verzí (což je libovolné \uv{slovo}, jediným účelem je
moci rozlišit mezi různými verzemi databáze téhož jména) a~typy uzlů, které obsahuje.
U~typů uzlů záleží na pořadí, v~němž jsou uvedeny. Důležité je, že definice indexů
na disk ukládána není, takže je možné používat jednu databázi v~různých programech
s~různými indexy. Každý index je určen jménem, pod nímž bude uživatelům přístupný,
svým typem a~typem uzlů, na který je navázán. Je-li jako typ uzlu uvedeno klíčové
slovo \verb|all|, je index informován o~změnách všech uzlů v~databázi.


\chapter{Ukázková implementace}

Tato kapitola popisuje specifika provedené implementace a~její omezení. Nejprve je
krátce pojednáno o~zvoleném programovacím jazyce a~platformě. Následuje popis
správy paměti a~dalších vnitřních mechanismů, především 
generování definic typů databází / uzlů
/ indexů uživatelem. Na závěr jsou zmíněny omezení provedené implementace
a~výsledky několika jednoduchých benchmarků.

Ukázková implementace je napsána
pod licencí GPL verze 2 \cite{GPLv2}. Jelikož nejsem zběhlý v~programování
v~assembleru a~překladač GCC neposkytuje vhodné atomické operace, převzal jsem
implementaci atomických operací z~linuxového jádra verze 3.3.4 \cite{linux3.3.4}.
(Jedná se o~soubor 
\verb|utils/atomic_amd64.h|, je pod licencí GPL verze 2, čímž implikuje
licenci pro celý projekt, neboť ten lze považovat za odvozené dílo.)


%%%%%%%%%%%%%%%%%%%%%%%%%%%%%%%%%%%%%%%%%%%%%%%%%%%%%%%%%%%%%%%%%%%%%%%%%%%%%%%%

\section{Jazyk implementace}

Ukázková implementace je provedena v~jazyce C, přesněji v~jeho verzi C99 \cite{C99}
s~použitím
GNU rozšíření a~vestavěných funkcí překladače GCC dostupných ve verzi 4.5
\cite{gccGNUexts}. Jako implementační platforma
byl zvolen Linux na architektuře AMD64.

Architektura AMD64
byla upřednostněna před IA-32,
protože na rozdíl od ní obsahuje všechny 8-bajtové
atomické instrukce (především \verb|xadd|, \verb|xchg| a~\verb|cmpxchg|)
nutné pro implementaci
verzovaných zámků a~dalších struktur pracujících atomicky (např. globální
hodiny či některé alokátory zmíněné níže vnitřně užívají atomické instrukce).
Implementaci je možné portovat i~na procesory architektury IA-32, které
podporují 8-bajtovou instrukci compare-exchange (\verb|cmpxchg8|), což jsou procesory
Intel Pentium a~novější. Je však nutné neexistující instrukce emulovat pomocí
\verb|cmpxchg8|.

Proč jsem použil GNU rozšíření jazyka C a~nezůstal u~standardizované verze C99 \cite{C99}?
Hlavním důvodem byla podpora operátoru \verb|typeof| a~vestavěných funkcí
\verb|__builtin_types_compatible| a~\verb|__builtin_choose_expr|. Díky nim je
možné implementovat makra \verb|trRead| a~\verb|trWrite| tak, že jako
parametry mají pouze uzel a~jméno jeho atributu s~nímž pracují, přičemž potřebná
typová informace se odvodí uvnitř makra.


%%%%%%%%%%%%%%%%%%%%%%%%%%%%%%%%%%%%%%%%%%%%%%%%%%%%%%%%%%%%%%%%%%%%%%%%%%%%%%%%

\section{Správa paměti}

\subsection{Opožděné uvolňování paměti}
Když transakce uvolní paměť, může nastat problém, pokud jiná transakce drží na tuto
paměť ukazatel a~pokusí se ho dereferencovat. Původní řešení v~TL2 enginu bylo
užití instrukcí, které při pokusu o~čtení paměti, kam nemá program přístup,
vrací nulovou hodnotu. Architektury IA-32 ani AMD64 žádné takové instrukce
nemají, a~tak se port TL2 enginu tyto architektury \cite{STMcur} s~tímto
problémem musel vypořádat jinak.

Řešením v~\cite{STMcur} bylo odchytávat signál neoprávněného přístupu do paměti.
Ač je tento signál možné odchytit, tak jedinou užitečnou informací, kterou
handler dostane je adresa, k~níž bylo přistupováno. Ta však nestačí k~ověření,
zda důvodem tohoto signálu byl skutečně přístup k~uvolněné transakční paměti
či chyba v~programu, takže handler musí signál prostě zahodit, čímž může
skrýt chybu v~programu. Dalším problémem je, že handler neoprávněného přístupu
do paměti se nesmí vrátit, neboť to způsobí nekonečnou smyčku. Proto musí být
užita nějaká varianta nelokálních skoků (např. volání \verb|setjmp| / \verb|longjmp|).
Z~důvodu režie
dává smysl použít \verb|setjmp| pouze na počátku transakce a~pomocí \verb|longjmp|
transakci rovnou restartovat, což je v~rozporu s~tím, že ukázková implementace
neobsahuje automatické restartování transakcí. Navíc registrací signal handleru
se můžeme dostat do konfliktu s~jinou částí aplikace, protože pro každý signál
může být registrován pouze jeden handler.
Naopak výhodou tohoto přístupu
je kromě jednoduchosti i~výkon, jelikož šance, že transakce $A$ načte ukazatel,
transakce $B$ uvolní paměť odkazovanou tímto ukazatelem a~transakce $A$ ukazatel
dereferencuje, je velmi malá.

V~této práci jsem místo tohoto přístupu použil opožděné uvolňování paměti --
paměť není uvolněna, když transakce žádá její uvolnění, ale až ve chvíli,
kdy je jisté, že žádná transakce na ni nemůže držet ukazatel.
Tohoto chování je docíleno tak, že transakční paměť se nealokuje pomocí volání
\verb|malloc|, ale pomocí volání \verb|node_allocator_alloc| pro alokaci uzlů
a~\verb|generic_allocator_alloc| pro zbylou paměť užívanou v~transakcích. Při uvolňování
je předán jako parametr i~čas, kdy byla paměť uvolněna. (Pokud byla paměť uvolněna
v~čase $x$, tak na ni mohou mít pointer pouze transakce, které začaly nejpozději
v~čase $x - 1$.) Skutečné uvolnění paměti ve správnou dobu pak zajistí příslušný
alokátor.

Výhodou opožděného uvolňování paměti je, že nemůže zakrývat chyby při
práci s~pamětí, ani nevyžaduje použití \verb|setjmp| / \verb|longjmp|. Nevýhodou
může být nutnost udržovat uvolněné bloky v~paměti a~s~tím spojená režie.
Tuto režii je možné vyvážit tím, že bloky k~uvolnění budeme používat
k~uspokojení požadavků na alokaci nové paměti, jak to činí \verb|VPageAllocator|.

\subsection{GenericAllocator}
\verb|GenericAllocator| je určen pro alokaci různě velkých bloků paměti užívané 
v~transakcích.
V~současné implementaci pouze zdržuje uvolnění paměti,
tím způsobem, že uvolňovanou paměť umístí do vnitřního seznamu a~když velikost
seznamu přesáhne danou mez (makro \verb|DB_GENERIC_ALLOCATOR_CACHE|),
tak zjistí počáteční čas nejstarší transakce $t$ a~uvolní ze seznamu paměť starší
než $t$. Uvnitř \verb|GenericAllocator|u je k~práci s~pamětí užívány funkce
\verb|malloc| a~\verb|free|, protože jejich implementace v~knihovně \verb|glibc|
je dobře optimalizovaná a~implementace rychlého alokátoru paměti je velmi
netriviální záležitost.

\subsection{NodeAllocator}
\verb|NodeAllocator| slouží k~alokaci uzlů. Zatímco \verb|GenericAllocator| je globální
pro každou instanci databáze, \verb|NodeAllocator| má v~rámci databáze každý typ
uzlu vlastní. Protože jsou všechny uzly jednoho typu stejně velké, je implementace
\verb|NodeAllocatoru| více optimalizovaná a~odpovídá struktuře alokátoru SLAB
\cite{BonwickSlabAllocator} pro malé objekty.
Tím je dána i~limitace současné implementace -- uzel nemůže být větší než cca 4000 bajtů
(velikost stránky minus hlavičky).

Jako v~alokátoru SLAB, jsou objekty alokovány z~větších bloků -- v~současné implementaci
bloky přímo odpovídají stránkám paměti -- a~ty jsou uloženy jako spojový seznam,
přičemž ne zcela zaplněné bloky jsou udržovány na konci seznamu. Je-li potřeba alokovat
prvek, je alokován z~bloku na konci seznamu. Pokud je tím blok zaplněn, je přesunut
na počátek seznamu. Pokud je poslední blok zcela plný, alokuje se nový blok.

Pro alokaci bloků je užíván \verb|VPageAllocator|. Vyprázdněný blok je okamžitě
předán \verb|VPageAllocator|u a~jako čas uvolnění je předán čas uvolnění posledního
uzlu z~daného bloku (což je korektní, protože posloupnost časů uvolnění je neklesající).

Významným důvodem pro použití
alokátoru typu SLAB byla potřeba šetřit místo.
Kromě ukazatelů spojového seznamu bylo možné z~uzlů odstranit i~ukazatel na
typ uzlu. Ten se přesunul na počátek bloku a~je k~němu možné se dostat díky tom,
že bloky odpovídají stránkám paměti, které jsou vždy zarovnané na svou velikost.
Pokud by uzly byly alokovány např. \verb|malloc|em, musel by v~hlavičce každého
z~nich být uložen ukazatel na typ uzlu, k~němuž přísluší, a~dva ukazatele, které
by všechny uzlu propojily do spojového seznamu (nutné pro jejich uvolnění
při uzavírání databáze a~pro vypsání databáze na disk). Připočteme-li režii
\verb|malloc|u, uspoříme použitím
alokátoru typu SLAB až 40 bytů na uzel, což
může být víc než je velikost užitečných dat v~daném uzlu.

\subsection{PageAllocator}
\verb|PageAllocator| slouží jako zdroj stránek pro jiné alokátory a~strukturu
\verb|FastStack| (implementaci zásobníku, která alokuje paměť po blocích, používanou
pro transakční log).
Stránky alokuje po blocích pomocí funkce \verb|mmap| a~nepoužité
stránky ukládá do spojového seznamu pro další využití. Pokud množství volných
stránek přesáhne daný limit, jsou navráceny systému voláním \verb|munmap|.
V~celé aplikaci typicky existuje jeden globální \verb|PageAllocator|.

\subsection{VPageAllocator}
\verb|VPageAllocator| je tenká vrstva
nad \verb|PageAllocator|em, která se stará o~opožděné
uvolňování paměti. Funguje jako \verb|GenericAllocator| s~tím
rozdílem, že stránky v~seznamu může použít pro uspokojení nových alokací.
Existuje jeden pro každou instanci databáze.


%%%%%%%%%%%%%%%%%%%%%%%%%%%%%%%%%%%%%%%%%%%%%%%%%%%%%%%%%%%%%%%%%%%%%%%%%%%%%%%%


\section{Omezení ukázkové implementace}
Protože implementace, která je součástí této práce, je zamýšlena spíše jako
ukázková
a~základ pro další testování různých implementací konkrétních
částí databáze (různé varianty servisního vlákna či přístupu k~zamykání atp.), než
jako kód připravený pro ostré nasazení, obsahuje také jistá omezení. Ta se pokusím
shrnout v~této sekci.

První omezení se týká práce se~signály. Pokud vlákno přijme signál a~začne vykonávat
uživatelem přiřazený handler ve chvíli, kdy je zablokováno v~určitých systémových
voláních (např. \verb|sem_wait|), skončí tato volání s~chybou
\verb|EINTR|, pokud vracejí chybový status. Co hůře systémové volání \verb|write|
může skončit
s~tím, že zapsalo pouze část dat. Chybě \verb|EINTR| se lze vyhnout, je-li při
registraci signal handleru použit příznak \verb|SA_RESTART|. S~ním systém dané volání
místo vrácení chyby restartuje. Protože současná implementace chybu \verb|EINTR|
neošetřuje, je při registraci handlerů signálů nutné použít příznak \verb|SA_RESTART|.

Další omezení spočívají v~implementaci alokátorů paměti. Méně významné z~nich souvisí
s~\verb|NodeAllocator|em. Protože ten staví na \verb|PageAllocator|u, je omezena
maximální možná velikost uzlu na cca 4000 bajtů (jedna stránka mínus hlavička přidaná
\verb|NodeAllocator|em). Vzhledem k~tomu, že největší z~atributů má velikost
10 bajtů (\verb|long double|), a~tedy uzel přesahující velikost 4000 bajtů by musel
mít přes 400 atributů, nepovažuji toto omezení za významné. Navíc k~jeho odstranění
by bylo potřeba alokaci paměti významně pozměnit (pravděpodobně by způsob alokace pro
uzly pod touto hranicí zůstal stejný a~byl by přidán speciální způsob alokace pro uzly
příliš velké (např. skrz \verb|GenericAllocator|)).

Poslední problém, který tu zmíním se týká \verb|GenericAllocator|u.
Současná implementace pouze odkládá uvolňování paměti (jak bylo popsáno výše), ale
není schopna paměť recyklovat. Hlavním důvodem k~tomuto chování je nemožnost zjistit
velikost u~bloku alokovaného pomocí funkce \verb|malloc|. Aby tedy bylo možné
bloky čekající na uvolnění použít znovu, bylo by nutné de facto implementovat
celý alokátor pro paměťové bloky různé velikosti. Jelikož toto samo o~sobě je
velmi komplikované, pokud chceme získat maximální výkon, a~s~ohledem k~optimalizacím,
které jsou v~GNU implementaci obsaženy (jelikož je dlouhodobě používána a~vylepšována),
jsem usoudil, že tato implementace bude stále rychlejší než neuměle napsaný obecný
alokátor s~integrovaným recyklováním bloků paměti.

\section{Výsledky testů}
Na závěr této kapitoly uvede ještě výsledky několika testů. Jde o~jednoduché syntetické
benchmarky, které prověřují škálování ukázkové implementace při provádění různých
operací a~při různých nastaveních databáze. Data, ze kterých byly vytvořeny tyto grafy,
naleznete v~příloze 2.

Všechny testy probíhaly na notebooku s~procesorem Intel Core i5 520M, který má
dvě fyzická jádra doplněná HyperThread technologií (takže má 4 logická jádra), 8~GB
paměti a~Linuxem verze 3.3.4. Frekvence procesoru byla pevně nastavena na 2133~MHz,
abych předešel zkreslení výsledků v~důsledku technologie TurboBoost.
Při všech testech byla databáze vytvořena s~příznakem \verb|DB_READ_ONLY|, který
způsobí, že výstup je místo do souboru s~databází zapisován do \verb|/dev/null|,
aby výsledky měření nebyly ovlivněny rychlostí pevného disku.

První test zjišťuje škálování uvedené implementace při čtení. V~každém testu byl
program spuštěn s~různým počtem vláken, přičemž mezi běhy zůstával stejný celkový
počet vykonaných transakcí a~odpovídajícím způsobem se měnil počet transakcí vykonaných
jedním vláknem. Z~grafu je patrné, že implementace dle očekávání škáluje do čtyř vláken
a~poté rychlost zůstává konstantní.

\centerline{\includegraphics[width=15cm]{mereni/read.pdf}}

Test zápisu také podává dobré výsledky, pokud spolu transakce nekolidují. Pokud
začnou pracovat na stejných datech, je zřejmé výrazné zpomalení v~důsledku restartování
transakcí. Jelikož cílovou skupinou jsou aplikace s~nízkým množstvím zápisů,
tak to příliš nevadí.

\centerline{\includegraphics[width=15cm]{mereni/write.pdf}}



\chapter*{Závěr}
\addcontentsline{toc}{chapter}{Závěr}

V~této práci jsem se pokusil navrhnout a~naprogramovat datovou strukturu, která
kombinuje výhody databází jako perzistenci a~transakce s~rychlostí běžných datových
struktur. Výsledkem práce je návrh
databáze, která udržuje data v~paměti, implementované jako knihovna
funkcí. Práce s~databází a~podpora transakcí je realizována pomocí konceptu
transakční paměti. Díky tomu je s~databází možné snadno pracovat paralelně z~více vláken.

Stále ale zbývá mnoho možností, jak současnou implementaci rozšířit či vylepšit.
Technicky nejjednodušší způsob zrychlení současné implementace je možnost přepsat
některé části, především výpočet CRC, s~pomocí SSE instrukcí a~tak je významně urychlit.
Mezi podstatně složitější, ale v~jistém směru užitečnější vylepšení patří implementace
mechanismu pro automatické restartování transakcí. 

Současná implementace sice obsahuje
makra \verb|trBegin| a~\verb|trCommit|, která automatické restartování umožňují,
ta ale používají
exponenciálně dlouhé čekání před dalším pokusem. To
je výhodné v~situaci, kdy jsou transakce malé a~množství kolizí nízké. Implementace
pokročilejšího algoritmu by nejspíše vyžadovala možnost nějak zvýhodnit konkrétní
transakce (především ty, které již opakovaně selhaly), což nyní není možné.

Dalším rozšířením, které by mohlo být velmi užitečné z~praktického hlediska, je zavedení
složených typů atributů jako jsou seznamy či slovníky (množiny párů klíč--hodnota
s~rychlým hledáním podle klíče). Hlavní obtíží v~tomto případě je návrh vhodného
rozhraní pro práci s~nimi a~serializaci daných operací na disk. Všechny současné atributy
si totiž vystačí s~operacemi \uv{přečti hodnotu} a~\uv{přiřaď hodnotu},
zatímco pro složené
atributy jsou tyto operace krajně nevhodné a~bylo by nutné zavést operace nové.

Posledním možným rozšířením, o~kterém se zmíním, jsou rozšířené kontexty indexů --
tedy že by kontext indexu nebyla pouze jeden objekt globální pro celou databázi, ale
indexy by mohly ukládat data i~do jednotlivých uzlů. To by mohlo zrychlit operace,
kdy je třeba uzel vyhledat v~rámci datové struktury představující index.



%%% Seznam použité literatury
%%% Seznam použité literatury je zpracován podle platných standardů. Povinnou citační
%%% normou pro bakalářskou práci je ISO 690. Jména časopisů lze uvádět zkráceně, ale jen
%%% v kodifikované podobě. Všechny použité zdroje a prameny musí být řádně citovány.

\def\bibname{Seznam použité literatury}
\begin{thebibliography}{99}
\addcontentsline{toc}{chapter}{\bibname}


\bibitem{HTM}
  {\sc Knight, Tom}.
  An architecture for mostly functional languages.
  In \emph{Proceedings of the 1986 ACM conference on LISP and functional programming},
  1986.

\bibitem{STM95}
  {\sc Shavit, Nir and Touitou, Dan}.
  Software Transactional Memory. 
  In \emph{Proceedings of the 14th ACM Symposium on Principles of Distributed Computing}, 
  pages 204–213. August 1995.

\bibitem{gccTM}
%  gcc TM FIXME
  Transactional Memory.
  In: \emph{GCC Wiki} [online, cit. 2012-05-23].
  Dostupné z: \url{http://gcc.gnu.org/wiki/TransactionalMemory}

\bibitem{haswellTM}
%  Haswell TM FIXME
  {\sc Reinders, James.}
  \emph{Transactional Synchronization in Haswell} [online], 2012 [cit. 2012-05-23].
  Dostupné z: \url{http://software.intel.com/en-us/blogs/2012/02/07/transactional-synchronization-in-haswell/}

\bibitem{STMcur}
  {\sc da Cunha, Goncalo Vasco Trincao Bento}.
  \emph{Consistent State Software Transactional Memory}. Master thesis.
  Faculdade de Ciencias e Tecnologia, Universidade Nova de Lisboa, 2007.

\bibitem{TL2engine}
  {\sc Dice, Dave, Shalev, Ori and Shavit, Nir}. Transactional Locking II. 
  In \emph{Proceedings of the 20th International Symposium on Distributed Computing} (DISC 2006),
  pages 194–208, Stockholm, Sweden, 2006.

\bibitem{ssd}
  %SSD zapisují po velkých blocích FIXME
  {\sc Shimpi, Anand Lal}.
  \emph{The SSD Anthology: Understanding SSDs and New Drives from OCZ} [online],
  2009 [cit. 2012-05-23].
  Dostupné z: \\ \url{http://www.anandtech.com/print/2738}

\bibitem{Redis}
  {\sc Sanfilippo, Salvatore et al.}
  \emph{Redis} [program], Version 2.4.14, 2009.
  %Key-value database Redis.
  Dostupné také z: \url{http://redis.io/download}

\bibitem{CRC32c}
  % CRC32c FIXME
  {\sc Koopman, Philip}.
  32-Bit Cyclic Redundancy Codes for Internet Applications. 
  In \emph{The International Conference on Dependable Systems and Networks}, 
  pages 459–468, 2002.

\bibitem{iSCSI}
  %iSCSI FIXME
  {\sc Satran, Julian et al.}
  RFC 3720, \emph{Internet Small Computer System Interface (iSCSI)}.
  The Internet Society, 2004.

\bibitem{btrfs}
%  btrfs FIXME
  %Main Page - btrfs Wiki.
  Btrfs design.
  In: \emph{btrfs Wiki} [online, cit. 2012-05-23].
  Dostupné z: \\ \url{https://btrfs.wiki.kernel.org/index.php/Btrfs_design}

\bibitem{sse42}
  %SSE 4.2 FIXME
  {\sc Intel Corporation}.
  \emph{Intel SSE4 Programming Reference} [online]. 2007 [cit. 2012-05-23].
  Dostupné z: \url{http://softwarecommunity.intel.com/isn/Downloads/Intel\%20SSE4\%20Programming\%20Reference.pdf}

\bibitem{utf8}
%  UTF-8 FIXME
  {\sc The Unicode Consortium}.
  \emph{The Unicode Standard}. Version 6.1.0.
  Mountain View, CA: The Unicode Consortium, 2012. ISBN 978-1-936213-02-3

\bibitem{IEEE754}
%  IEEE 754. FIXME
  IEEE 754:2008, \emph{IEEE Standard for Floating-Point Arithmetic}.
  The Institute of Electrical and Electronics Engineers, Inc., 2008.

\bibitem{linuxCOW}
 % Linux COW FIXME
  {\sc Love, Robert}.
  \emph{Linux System Programming: Talking Directly to the Kernel and C Library}.
  O'Reilly Media, Inc., 2007.

\bibitem{Perl}
%  Perl FIXME
  {\sc Wall, Larry et al.}
  \emph{Programming Language Perl} [program]. Version 5.4.12. 2010.
  Dostupné také z: \url{http://www.perl.org/get.html}

\bibitem{RecDescent}
%  RecDescent FIXME
  {\sc Conway, Damian and Braun, Jeremy T.}
  \emph{Parse::RecDescent -- Generate Recursive-Descent Parsers} [program].
  Version 1.965.1. 2003. % [cit. 2012-05-23].
  Dostupné také z: \url{http://search.cpan.org/~dconway/Parse-RecDescent-1.965001/lib/Parse/RecDescent.pm}

\bibitem{GPLv2}
  %GPL v2 FIXME
  {\sc Free Software Foundation, Inc.}
  \emph{GNU General Public License}. Version~2. 1991.

\bibitem{linux3.3.4}
  % Linux 3.3.4 FIXME
  {\sc Torvalds, Linus et al.}
  \emph{Linux} [program]. Version 3.3.4. 2012.
  Dostupné také z: \url{http://www.kernel.org/pub/linux/kernel/v3.x/linux-3.3.4.tar.gz}

\bibitem{C99}
%  C99 FIXME
  ISO/IEC 9899:1999, \emph{Programming Language C}.
  International Organization for Standardization 1999. 

\bibitem{gccGNUexts}
%  GCC GNU exts. FIXME
  {\sc Free Software Foundation, Inc.}
  C Extensions. \emph{Using the GNU Compiler Collection (GCC)} [online]. Version 4.5.0.
  2008 [cit. 2012-05-23].
%  C Extensions - Using the GNU Compiler Collection (GCC).
  Dostupné z: \url{http://gcc.gnu.org/onlinedocs/gcc-4.5.0/gcc/C-Extensions.html}

\bibitem{BonwickSlabAllocator}
  {\sc Bonwick, Jeff.}
  \emph{The Slab Allocator: An Object-Caching Kernel Memory Allocator}.
  Sun Microsystems, 1994.
  % FIXME

\end{thebibliography}



\chapwithtoc{Příloha 1: Ukázkový program}
Jednoduchý program využívající hybridní databázi. Obsahuje databázi
jmen a~pozdravů. Umožňuje přidat jméno či pozdrav a~vypsat náhodný pozdrav.

\lstinputlisting[language=C,
basicstyle=\footnotesize\tt,           % the size of the fonts that are used for the code
  numbers=left,                   % where to put the line-numbers
  stepnumber=1,                   % the step between two line-numbers. If it's 1, each line 
                                  % will be numbered
  numbersep=5pt,                  % how far the line-numbers are from the code
  showspaces=false,               % show spaces adding particular underscores
  showstringspaces=false,         % underline spaces within strings
  showtabs=false,                 % show tabs within strings adding particular underscores
  tabsize=2,                      % sets default tabsize to 2 spaces
  captionpos=b,                   % sets the caption-position to bottom
  breaklines=true,                % sets automatic line breaking
  breakatwhitespace=true]{priklad.c}


\chapwithtoc{Příloha 2: Výsledky testů}

\section*{Měření rychlosti čtení}
Hodnoty v~tabulce udavájí čas v~milisekundách, který trvalo provedení
milionu transakcí. Za proměnnou \verb|$i| byl dosazen počet vláken.
Testovací databáze byla vytvořena příkazem \verb|create2.exe| a~obsahovala
milion prvků.

\setlength{\parindent}{0cm}
Příkaz: \verb|./read.exe /tmp/test_db $i $(((1000 * 1000) / $i)) 50 1|
\setlength{\parindent}{1cm}

\begin{center}
  \input{mereni/read.tex}
\end{center}

\section*{Měření rychlosti zápisu}
Hodnoty v~tabulce udavájí čas v~milisekundách, který trvalo provedení
200~000 transakcí. Za proměnnou \verb|$i| byl dosazen počet vláken.
Testy s~kolizemi a~bez kolizí se liší pouze tím, zda si pracovní vlákna
rozdělila databázi na části a~každé pracovalo se svou částí nebo zda
všechna vlákna přistupovala k~celé databázi.
Testovací databáze byla vytvořena příkazem \verb|create2.exe| a~obsahovala
100~000 prvků.


\subsubsection*{Bez kolizí:}
\setlength{\parindent}{0cm}
Příkaz: \verb|./write.exe /tmp/test_db $i $(((1000 * 200) / $i)) 50 1|

\begin{center}
  \input{mereni/write.tex}
\end{center}

\newpage

\subsubsection*{S kolizemi:}
Příkaz: \verb|./write_col.exe /tmp/test_db $i $(((1000 * 200) / $i)) 50 1|

\begin{center}
  \input{mereni/write_col.tex}
\end{center}

\setlength{\parindent}{1cm}
\subsubsection*{Kolize:}
Následující tabulka obsahuje počet kolizí vyjádřený jako promile úspěšně dokončených
transakcí. Data pocházejí z~běhu zápisového testu s~kolizemi.

\begin{center}
  \input{mereni/col.tex}
\end{center}




%%% Tabulky v bakalářské práci, existují-li.
% \chapwithtoc{Seznam tabulek}

%%% Použité zkratky v bakalářské práci, existují-li, včetně jejich vysvětlení.
% \chapwithtoc{Seznam použitých zkratek}

%%% Přílohy k bakalářské práci, existují-li (různé dodatky jako výpisy programů,
%%% diagramy apod.). Každá příloha musí být alespoň jednou odkazována z vlastního
%%% textu práce. Přílohy se číslují.
% \chapwithtoc{Přílohy}

\openright
\end{document}
